% SOURCE_PDF_PAGE: 053
\chapter{Gordle 即服务}

\begin{infobox}{本章内容}
\begin{itemize}
  \item 创建并运行一个 HTTP 服务器，在给定端口上监听消息
  \item 用 GET 和 POST 动词监听端点
  \item 构建带状态码的响应
  \item 解码不同来源的数据：路径与查询参数、请求体与请求头
  \item 使用正则表达式进行测试
\end{itemize}
\end{infobox}

1962 年，J.\,C.\,R.~Licklider 提出了让计算机通过网络相互通信的可能。从那时起，计算机科学走过了漫长的道路：先是穿过他命名的``星系际计算机网络''（Intergalactic Computer Network），然后是先驱性的高级研究计划局网络（ARPANET），再到今天的互联网。如今网络已融入日常——当你拿起手机查天气、看新闻、甚至只是看时间时，都在用它。能够从远程位置使用服务器这件事至关重要，尤其是在 2020 年，整个星球都因新冠疫情进入封锁状态的时候。

说到底，服务器其实就是一台机器，它在给定的一组端口上监听通信，并能够应答收到的消息。本章我们将实现这样一个服务器。为了让事情有趣一点，我们的服务器将提供一个 API，让用户可以玩一局 Gordle。本章我们聚焦于服务器端的实现，而不是算法层面——那已经在第 5 章讲过了。

% SOURCE_PDF_PAGE: 054
本章第一部分，我们将创建 REST API，并用简单的工具测试它，比如浏览器或一条命令。第二部分，我们将集成 Gordle 游戏——这需要对它做一些更新，以符合我们在服务器中的使用方式。最后，我们会提几点安全方面的建议。

\begin{infobox}{注意}
容器化与部署等后续步骤不在本章的话题范围内。更多信息请参阅 Marko Lukša 与 Kevin Conner 的《Kubernetes in Action》（Manning，2017，\url{www.manning.com/books/kubernetes-in-action}）。
\end{infobox}

本项目有以下要求：

\begin{itemize}
  \item 在一个 web 服务上玩 Gordle。
  \item 运行一个至少暴露以下端点的服务：
  \begin{itemize}
    \item \texttt{NewGame} 创建一个会话，返回一个玩家 ID。它将用于统计剩余尝试次数。
    \item \texttt{GetStatus} 返回——嗯——一局游戏的状态，包括还允许多少次猜测、之前用过的提示，等等。
    \item \texttt{Guess} 接收一个单词作为参数，返回反馈和游戏状态。
  \end{itemize}
  \item 第二步中，加入对玩家会话的某种跟踪。许多在线资源都使用某种用户标识——多数时候是身份验证。你会看到我们如何传递这一信息。
\end{itemize}

以下是本项目不要求的：监控、日志策略，以及可扩展性（暂时）。

% SOURCE_PDF_PAGE: 055
\section{新服务的空壳}

可以把 web 服务想象成一个守护进程（daemon）——它一直在运行，等待发送到其所宿主计算机某个端口上的查询。本章我们将采用由外而内的方式开发：先创建一个监听端口但什么都不做的服务；然后添加空的端点，每个功能一个；最后再填入它们的逻辑。

当团队里的其他人——例如前端开发者，或者干脆是其他团队——正等着你的成果时，这种策略最为合适。你可以先返回一个静态的模拟响应，让别人先用起来，并在你开发的同时给你反馈。一个能返回点什么的服务——哪怕返回的是常量或不相干的内容——往往就足以帮助别人设计、开发、测试或部署他们的方案。在开始写几行代码之前，先介绍一些词汇并理解几个理论概念是很重要的。

\subsection{服务器、服务、web 服务、端点与 HTTP 处理器}

服务器可以想象成一台运行在某个地方的计算机。我们通常让服务器一直运行，因为它们承载着服务（service），即暴露 API 的应用程序。通常我们希望服务永久运行，因为停下来的服务毫无用处。web 服务是一种特殊的服务，它利用 web 协议来接收和发送与外界的通信。

\begin{infobox}{注意}
记住本章后面的步骤：服务不应终止自己的执行。换句话说，它会一直运行，直到用户决定停掉它。
\end{infobox}

% SOURCE_PDF_PAGE: 056
端点（endpoint）是外部进入服务的访问点。对 web 服务而言，端点映射到特定的 URL，你很快就会看到。web 服务在幕后会使用一个 HTTP 处理器（HTTP handler）来处理请求。HTTP 处理器接收 HTTP 请求并生成 HTTP 响应。

\begin{infobox}{注意}
请注意，web 服务和端点这两个术语在不同语境下的定义略有差别，但在本章的讨论中，我们将使用刚才给出的定义。
\end{infobox}

\subsection{动手写代码}

先为我们的服务创建模块。记着我们要在 HTTP 服务器里运行 Gordle，就能起一个贴切的名字。记住，如果你打算把代码推送到代码仓库，最好把仓库的完整路径声明为模块名：

\begin{shellcode}[]
> go mod init learngo/httpgordle
\end{shellcode}

创建好 \texttt{go.mod} 文件后，就可以开始写 \texttt{main} 函数了，它将负责创建并运行一个服务器（见代码清单 8.1）。为此，我们需要 \texttt{net/http} 包帮点忙——第 6 章已经见过它。它的文档相当长，不过 \texttt{go doc net/http} 的头几行就说明：``http 包提供了 HTTP 客户端与服务器的实现。''目前我们只关心服务器那一半。

服务器要监听一个特定端口，所以在你的主机上随便找个空闲端口。HTTP 的默认端口是 80，但出于开发目的，我们更愿意用别的端口，比如 8000 或 8080，因为 80 在你的机器上
% SOURCE_PDF_PAGE: 057
多半已被别的东西占用。\texttt{net/http} 包里有一个函数看起来恰好能满足需求——\texttt{ListenAndServe}。

\begin{gocode}[title={代码清单 8.1 \texttt{main.go}：创建服务器}]
package main

import "net/http"

func main() {
    // Start the server.
    err := http.ListenAndServe(":8080", nil) #1
    if err != nil {
        panic(err) #2
    }
}

#1 Use any free port here.
#2 It's OK to panic in the body of main.
\end{gocode}

\begin{codeannotations}
\item 这里用任何空闲端口都行
\item 在 main 函数体里 panic 是可以的
\end{codeannotations}

希望这没有吓到你。我们给 \texttt{ListenAndServe} 传了两个参数。第一个是我们希望 web 服务监听的地址——\texttt{localhost:8080}（\texttt{localhost} 可省略）是个热门选择。如果你在 Windows 上，考虑用 8080 以外的端口，以免弹出 Windows 防火墙提示。第二个参数是能处理请求的处理器（下一节讨论）——现在可以先传 \texttt{nil}，但逻辑将来就要实现在这里。我们说过在 main 函数体里 panic 没问题，但其实这只对了一半。\texttt{panic} 会倾倒整个堆栈跟踪，这会让用户摸不着头脑。更礼貌的做法是输出错误并以合适的退出码结束，比如下面这段代码（代价是无法保证所有 defer 调用都已执行）：

\begin{gocode}[]
if err != nil {
    fmt.Fprintln(os.Stderr, err)
    os.Exit(1)
}
\end{gocode}

跑起来看看：

\begin{shellcode}[]
> go run .
\end{shellcode}

% SOURCE_PDF_PAGE: 058
你可能注意到执行``卡住''了。这是因为 \texttt{ListenAndServe} 函数永不返回。毕竟，这正是我们想要的行为：服务正在运行！

手动测试 HTTP 服务器的工具很多。这里列出四种：

\begin{itemize}
  \item 互联网浏览器——浏览器生来就是发 HTTP 请求的。设置某些参数（比如请求体、请求头或想用的动词）可能有点麻烦，但对这第一个实现来说够用了。
  \item Postman——这个工具有图形界面，在网络上传送格式化消息时比 web 浏览器更精细。
  \item \texttt{curl}——这个命令行工具暴露了我们需要的一切。这是我们的首选：不仅因为在书里分享命令行更简单，还因为使用命令行工具而非图形界面会让测试更容易自动化。\texttt{curl} 随每个版本的 Linux 或 macOS 附带，Windows 10 及以上版本也有。
  \item Go 程序——你可以写一个 Go 程序，创建一个客户端与我们的服务器对话，就像第 6 章那样。
\end{itemize}

我们提供的 nil 处理器其实做不了什么，但还是可以看看它的实际表现。打开你最喜欢的浏览器，输入 URL \texttt{localhost:8080}，你会看到默认 HTTP 处理器以一条 404 消息的形式给出的响应。

如果你想用 \texttt{curl}，同样的请求写成命令行就是：

\begin{shellcode}[]
> curl http://localhost:8080
\end{shellcode}

% SOURCE_PDF_PAGE: 059
它同样会返回一条 404 消息。在接下来的测试中我们会更多地使用 \texttt{curl}。

第一个 HTTP 服务器成功实现了。在进入下一节之前，需要先杀掉我们的服务器。在执行 \texttt{go run .} 的那个 shell 里按 Ctrl–C。这会给当前进程发送一个中断信号，将其终止。终止运行中的程序还有其他办法：用你电脑的任务管理器，或者使用 \texttt{kill} 命令——如果你知道想终止的程序的进程 ID。做完这些后，先提交一下工作，再进入下一部分——处理掉这个 nil 处理器。

\section{添加端点}

本项目要实现的是一个让人们玩 Gordle 的 web 服务。再看看需求：我们需要能创建一局游戏、进行一次猜测、获取一局游戏的状态。我们立刻可以看出，我们的服务要打交道的实体是``游戏''——创建它们、使用它们、展示它们。我们还需要能访问这些游戏——为了猜一次或查看状态。为此，我们需要某种标识游戏并存储它的方式。存储稍后再说，因为我们是由外而内开发的。

HTTP 服务器上的端点是一对``路径 + HTTP 方法''。路径应反映正在使用哪个资源。在我们的例子里，打交道的是游戏，路径 \texttt{/games} 是个好的起点。需要标识单个游戏时，可以用 \texttt{/games/\{game\_id\}}。

HTTP 方法（或动词）描述调用地址时想执行的动作。已定义的方法有好几种，但我们聚焦于下面这些：

% SOURCE_PDF_PAGE: 060
\begin{description}
  \item[GET] 用于访问资源。
  \item[POST] 用于创建资源，或请求数据被处理。
  \item[PUT] 用于更新资源。
  \item[DELETE] 用于删除资源。
\end{description}

其中一些端点——在这个列表里是 GET、PUT 和 DELETE——应当是幂等的；也就是说，连续多次调用都应返回相同的响应，并让资源处于与只调用一次时相同的状态。如果你的 API 不幂等，就需要在文档里明确告诉用户。GET、PUT 和 DELETE 干脆不应被用于非幂等的端点——那种场合请改用 POST。现在，我们可以实现第一个端点了。

\subsection{创建新游戏}

创建游戏是玩家要做的第一件事。先想想我们要达成的目标。创建一局 Gordle 其实不需要任何输入。如果你还记得第 5 章，我们当年是用 \texttt{go run main.go} 启动游戏的。现在要新增一个端点，就需要一对``路径 + HTTP 方法''。我们要使用的资源是游戏——\texttt{/games} 这个路径看起来完美。

该用哪个 HTTP 方法？本例中我们在创建游戏，所以应该用 POST。某些情况下，你可能已经知道要创建资源的标识符。例如，处理图书时，我们可以用 ISBN 配合 PUT 方法在地址 \texttt{/books/9781633438804} 上创建图书资源。对 Gordle 而言，我们只需要创建一个空游戏。

不过，为了能跟踪它，至关重要的一点是返回游戏的标识符，\texttt{GetStatus} 和
% SOURCE_PDF_PAGE: 061
\texttt{Guess} 端点都要用到它。标识符可以是一串数字（比如电话号码就是纯数字标识符），也可以是字符（比如车牌就是一辆车的标识符）。有一些不错的库能提供唯一标识符，比如 \url{https://github.com/google/uuid}。我们将坚持用随机整数，因为这个话题在第 5 章需要从列表中随机取词时已经讲过了。

这个端点的 API 就定义好了。我们知道了路径、与之关联的动词，以及期望的响应——一个标识符。文档可以这样开头：

\begin{gocode}[]
POST /games - creates a new game and returns its ID.
\end{gocode}

好，回到代码！

\subsubsection*{项目组织}

关于模块内的文件如何组织、项目内的模块如何组织，并没有官方规定。不过，有些常见实践值得一提。第一个要点是，有几个文件夹名字在 Go 中有特定行为：

\begin{description}
  \item[testdata] 第 3 章已经提过，名为 \texttt{testdata} 的目录不会被 Go 工具检视。这些目录里的代码不会被 \texttt{go build} 或 \texttt{go run} 编译，里面写的测试不会被 \texttt{go test} 执行，文档也不会通过 \texttt{go doc} 展示。
  \item[internal] 现在是介绍另一个特殊目录名的时候了：\texttt{internal}。internal 目录可以包含供当前模块使用的代码，但这些代码对其他模块不可见。例如，模块 \texttt{golang.org/x/text} 就有一个 internal 包，
% SOURCE_PDF_PAGE: 062
其中定义了 \texttt{InheritanceMatcher} 类型。然而，尽管这个类型是导出的（因为首字母大写），我们还是无法在自己的模块里创建这种类型的变量。定义在名为 internal 的目录（或 internal 目录的子目录）中的类型和函数，其作用域仅限于当前模块（在我们举的例子里，即 \texttt{golang.org/x/text} 模块）。internal 目录是安放``你不想让别人用的代码''的好地方。对一个服务来说，大部分代码都会住在那里。

\item[vendor] 提 \texttt{vendor} 目录纯粹是出于历史原因。我们不打算细讲这门语言的整部历史，但早期版本的 Go 曾经使用``带版本''的依赖——每个模块本地的副本。这些副本会被放进 vendor 目录——在你的版本控制工具里忽略该目录的内容曾是个好主意。出于兼容性考虑，最好干脆永远不要把目录命名为 vendor。实在要用，就叫 \texttt{vendors} 吧。

\item[pkg] 你可能会遇到位于模块根目录 \texttt{pkg} 包里的包：\texttt{module/pkg/my\_package}。在 \texttt{pkg} 里，你通常能找到可以在你的项目之外使用的库。我们不鼓励使用 \texttt{pkg}，因为它不是 Go 的标准。它更像一个历史遗留物，或者说是 \texttt{golang-standards/project-layout} 的产物——那并不是 Go 团队的官方标准。
\end{description}

以上是硬性规则，接下来是可以自由跟随的建议。我们喜欢把服务的 API 放在一个 \texttt{api} 包里——模块根目录下的一个目录。

% SOURCE_PDF_PAGE: 063
\subsubsection*{服务的文件组织}

现在可以整理代码了：在模块根目录创建一个 \texttt{api} 目录，再创建一个 \texttt{internal} 目录，我们会在其中写各种各样的东西，包括放在子目录里的 HTTP 处理器。端点的实现方式对外部开发者毫无用处，所以干脆把它藏进 \texttt{internal/handlers} 目录。

我们选择为每个处理器创建一个包。就本项目而言，把所有处理器放进一个简单的包也完全没问题；不过，通过使用多个包，我们可以演示如何组织一个更大的项目，而不用真去写一个大项目。视情况而定，你有时可以在同一个包里完成所有事情（为了清晰起见分成不同文件），也可以把逻辑分散到不同的包。这里，我们选择把逻辑（校验、调用存储等）留在处理器里（它们负责与 HTTP API 相关的一切），因此我们偏好每个端点一个包。跟往常一样，做这类决策时，要想想增加功能后每个包将如何扩展和生长。

最后，很重要的一点是考虑耦合的概念。有些结构和函数天然耦合，必须一起更新。耦合紧密的代码分散到越多的包里，维持代码质量就越困难。跟往常一样，你需要找到平衡点；如果日后的组织方式不再符合需求，后续更新自然会提供重构组织的机会。

我们服务的 HTTP API 可以放在 \texttt{http.go} 文件中，而新游戏的初始处理器将放在 \texttt{newgame/handler.go} 文件里。我们将在 \texttt{router.go} 文件中把 API 绑定到处理器。此时的文件组织如下：

% SOURCE_PDF_PAGE: 064
\begin{shellcode}[]
> tree
.
├── go.mod
├── internal
│   ├── api
│   │   ├── doc.go
│   │   └── http.go
│   └── handlers
│       ├── doc.go
│       ├── newgame
│       │   └── handler.go
│       └── router.go
└── main.go
\end{shellcode}

\subsubsection*{定义 REST API}

REST 是一套约定，帮助我们在 HTTP 服务器上定义 API。它定义了资源的集合（collection）以及与它们交互的方式。

先从 \texttt{http.go} 文件开始。它应当包含我们需要暴露的一切，让别人能够使用我们接下来要实现的 \texttt{NewGame} 端点。所谓``一切''，指的是应当使用哪个 URL、哪个方法。如果还有别的，比如查询参数，也一并放进这个文件。在 \texttt{internal/api} 包中创建 \texttt{http.go} 文件。

\begin{gocode}[title={代码清单 8.2 \texttt{api/http.go}：定义 \texttt{NewGame} 端点所需之物}]
package api

import "net/http"

const (
    // NewGameRoute is the path to create a new game.
    NewGameRoute  = "/games"
)
\end{gocode}

现在，我们有了即将实现的 \texttt{NewGame} 端点所需的常量。这个端点的调用者应当期待什么返回内容？有些创建端点只返回一个 ID。这里我们想更详细一些，返回我们创建的完整游戏：我们 Gordle 游戏的客户端需要知道
% SOURCE_PDF_PAGE: 065
秘密单词的字符数和允许的最大尝试次数。既然我们要在 API 里定义``游戏''是什么，就应当考虑想包含的每一个字段。还可以告诉对方游戏的状态，让玩家知道自己能否继续玩、是已经赢了还是输了。最后，附上之前尝试的列表有助于展示。我们来把 JSON 游戏的形状定义如下：

\begin{gocode}[]
{
  "id": "1225482481867118141",
  "attempts_left": 4,
  "word_length": 5,
  "status": "Playing",
  "guesses": [
    {"word":"slice","feedback":""}
  ],
}
\end{gocode}

这能轻松翻译成一个带 JSON 标签的 Go 结构体，如前几章所示，也如下面的代码清单所示。

\begin{gocode}[title={代码清单 8.3 \texttt{internal/api/http.go}：定义游戏的 API 结构}]
package api

// ...

// GameResponse contains the information about a game.
type GameResponse struct {
    ID           string  `json:"id"`
    AttemptsLeft byte    `json:"attempts_left"`
    Guesses      []Guess `json:"guesses"` #1
    WordLength   byte    `json:"word_length"`
    Solution     string  `json:"solution,omitempty"` #2
    Status       string  `json:"status"`
}

// Guess is a pair of a word (submitted by the player)
// and its feedback (provided by Gordle).
type Guess struct {
    Word     string `json:"word"`
    Feedback string `json:"feedback"`
}

#1 Guesses come as a list of previous attempts.
#2 The solution is only given when the game is over; it's omitted
otherwise.
\end{gocode}

\begin{codeannotations}
\item \texttt{Guesses} 以之前尝试列表的形式出现
\item 只有游戏结束时才给出答案；否则省略
\end{codeannotations}

% SOURCE_PDF_PAGE: 066
注意，所有类型都是基本类型；我们没有对使用方强加任何强类型。我们选择用字符串来表达反馈。等开始填充反馈时再细说。现在结构有了，也正式发布给使用方了，是时候让服务器真正暴露这个端点了。

\subsubsection*{HTTP 多路复用器、Handle 与 Handler}

如果你还记得第一节，我们当时给 \texttt{ListenAndServe} 函数传了个 nil 处理器。仔细看看这个函数的第二个参数，它是一个 \texttt{http.Handler}。该类型声明如下：

\begin{gocode}[]
type Handler interface {
    ServeHTTP(ResponseWriter, *Request)
}
\end{gocode}

实现一个可以传给 \texttt{ListenAndServe} 的 Handler 有好几种方式。比如，我们可以定义一个结构体，让它实现这个接口：

\begin{gocode}[]
type newGameHandler struct {}

func (h *newGameHandler) ServeHTTP(
    w *http.ResponseWriter, req *http.Request) {
    ...
}
\end{gocode}

然后在 main 函数里这样用：

\begin{gocode}[]
err := http.ListenAndServe(":8080", newGameHandler{})
\end{gocode}

这样也行。但这里有个小问题：这样创建处理器只能定义一个端点。而我们知道，我们要实现好几个端点——至少三个。

再看看 \texttt{net/http} 包提供的另一个类型：\texttt{ServeMux}。一条快速的 \texttt{go doc http.ServeMux} 命令
% SOURCE_PDF_PAGE: 067
就能看出它是一个请求多路复用器（multiplexer）。也就是说，\texttt{ServeMux} 负责根据请求被发送到的 URL 来分发请求。多路复用器（常称 mux）是 HTTP 服务的基石。

关于 \texttt{ServeMux} 还有两个要点值得强调：第一，它允许用 \texttt{HandleFunc} 方法注册端点，这正是我们想做的。第二，\texttt{ServeMux} 有一个签名正确、名为 \texttt{ServeHTTP} 的方法。它实现了 \texttt{Handler} 接口，因此我们可以把一个 \texttt{ServeMux} 传给 \texttt{ListenAndServe} 函数。

\subsubsection*{多路复用器}

先来写多路复用器，然后再看它注册端点时接受什么签名。

我们需要构建一个 \texttt{http.Handler}。起初只要几行，但服务一旦成长，它会长得飞快。所以我们写一个函数来构建并返回它，如代码清单 8.4 所示。构建 mux 的全部逻辑都隔离在这里。

这个函数创建一个新实例，让它监听我们未来的端点，然后返回。我们还没定义端点长什么样，所以先放一个占位符。如果你想编译一下、检查其他部分是否成立，\texttt{nil} 完全可以接受；不过如果用 \texttt{nil}，别指望对你的服务发起请求会有什么结果——除了 panic。

% SOURCE_PDF_PAGE: 068
\begin{gocode}[title={代码清单 8.4 \texttt{internal/handlers/router.go}：处理器与 URL 的关联}]
package handlers

import (
    "net/http"

    "learngo-pockets/httpgordle/internal/api"
    "learngo-pockets/httpgordle/internal/newgame"
)

// Mux creates a multiplexer with all the endpoints for our service.
func Mux() *http.ServeMux {
    mux := http.NewServeMux()
    mux.HandleFunc(api.NewGameRoute, newgame.Handle) #1
    return mux
}

#1 Connects a URL to a handler
\end{gocode}

\begin{codeannotations}
\item 把 URL 连接到处理器
\end{codeannotations}

我们终于可以在 main 函数里使用这个 \texttt{Mux} 了，替换之前的 nil 处理器。

\begin{gocode}[title={代码清单 8.5 \texttt{main.go}：使用新的 \texttt{Mux()} 函数}]
package main

import (
    "net/http"

    "learngo-pockets/httpgordle/internal/handlers"
)

func main() {
    err := http.ListenAndServe(":8080", handlers.Mux())
    if err != nil {
        panic(err)
    }
}
\end{gocode}

现在，剩下的就是实现我们传入 mux 的那个 \texttt{newgame.Handle} 处理器。接下来就做这件事。

\subsubsection*{NEWGAME \{ENDPOINT\} 的处理器}

\texttt{mux.HandleFunc} 方法的签名如下：

\begin{gocode}[]
func (mux *ServeMux) HandleFunc(
    pattern string, handler func(ResponseWriter, *Request))
\end{gocode}

% SOURCE_PDF_PAGE: 069
这个方法为给定的路径注册一个处理器——也就是我们作为第二个参数传入的匿名函数。相比 \texttt{http.Handle}，这个方法的好处是我们不必编写新的 \texttt{http.Handler}；只需提供处理器本身，即负责处理请求并写出响应的那个函数。

在最简单的实现里，我们可以在 \texttt{ResponseWriter} 上调用 \texttt{Write} 函数，对客户端说点什么。把处理器写到 \texttt{internal/handlers/newgame/handler.go} 文件里。

\begin{gocode}[title={代码清单 8.6 \texttt{handler.go}：空的 newgame 处理器}]
package newgame

import "net/http"

func Handle(w http.ResponseWriter, req *http.Request) {
    _, _ = w.Write([]byte("Creating a new game"))
}
\end{gocode}

这是最初的版本。我们连错误都还没检查，不过随着逐步接近最终版本，会加上检查的。

好，越来越像样了！现在可以 \texttt{go run .} 运行这段代码，看看它的表现。在进入下一节之前，还有最后一步要做。让我们一起发现它。

如果服务正在运行，你在浏览器打开 \texttt{localhost:8080/games}，或者运行 \texttt{curl http://localhost:8080/games}，你应该会收到一条消息，告诉你正在创建一局游戏。这些工具传递了一个不可见的参数——我们没有指定用哪个 HTTP 方法，两个工具都默认选择并发送了 GET。我们希望 \texttt{NewGame} 端点只接受 POST 请求。来实现这一点。

% SOURCE_PDF_PAGE: 070
\subsubsection*{请求会被收到}

虽然听起来像是废话，但一个服务器应当能同时服务多个客户端。某人正在使用一个端点，不应妨碍其他人也来调用它。在幕后，这意味着服务器不会等一个任务完成才去服务新的调用。在 Go 中，这是用 goroutine 实现的。goroutine 是 Go 版本的并发编程——其他语言里最接近的等价物通常叫作线程、协程、纤程或绿色线程。不过 goroutine 与线程并不相同，我们会在第 9 章更全面地介绍 goroutine。

当服务器收到一个请求，它会启动一个 goroutine 去执行处理器。这看起来美妙又方便，但也有其局限。本章最后一节讲正确性，会呈现两个重要话题——竞态条件，以及确保服务器在试图同时处理两个请求时不会崩掉。如第 7 章所述，我们不运行 \texttt{go test .}，而是加上 \texttt{-race} 标志。

\subsubsection*{HTTP 状态码与请求头}

一局 Gordle 只应在收到 POST 请求时创建。但如果收到了错误的方法该怎么办？又该如何实现？

第一个问题的答案很明确：我们应当拒绝这条消息。其实 HTTP 有一个专门给错误动词准备的状态码，那就不妨用它：\texttt{http.StatusMethodNotAllowed}（见表 8.1）。至于在哪里做这个检查，唯一合乎逻辑的位置就是 \texttt{Handle} 函数内部。

% SOURCE_PDF_PAGE: 071
HTTP 响应头通过调用 \texttt{WriteHeader} 并传入合适的状态码来设置。如果收到的请求不是我们在 API 里声明的那个方法，我们可以在那里立即结束这次调用。下面的代码清单展示了在 \texttt{internal/handlers/newgame/handler.go} 文件的 \texttt{newGameHandler} 函数中加入检查的实现。

\begin{gocode}[title={代码清单 8.7 \texttt{handler.go}：拒绝使用错误方法的请求}]
func Handle(w http.ResponseWriter, req *http.Request) {
    if req.Method != http.MethodPost {
        w.WriteHeader(http.StatusMethodNotAllowed)
        return
    }
    _, _ = w.Write([]byte("Creating a new game"))
}
\end{gocode}

可以看到，这与``常规''Go 代码里的错误处理并不相同。我们必须适应 HTTP，这意味着要通过状态码来传达错误。既然处理一个无效的请求毫无意义，我们可以放心地从处理器返回。错误——无效的方法——已经处理完毕，我们没有别的事可做了。

运行之前的测试：启动服务，用各种工具访问 \texttt{http://localhost:8080/games} 页面。视浏览器而定， 浏览器的不同，你可能会、也可能不会看到 405 错误。在我们的测试里，Firefox 什么都没显示，而 Google Chrome 显示了。用 \texttt{curl} 看看：

\begin{shellcode}[]
> curl localhost:8080/games
\end{shellcode}

什么都不返回——一个完全空的响应。这就有问题了，因为它让我们更难检查实现。所幸 \texttt{curl} 也自带选项，其中一个有意思的是 \texttt{-\-verbose}（或 \texttt{-v}），它会打印多得多的信息。试一试吧——你应该会得到类似这样的输出：

% SOURCE_PDF_PAGE: 072
\begin{shellcode}[]
> curl -v localhost:8080/games
    Output:
*   Trying 127.0.0.1:8080...
* Connected to localhost (127.0.0.1) port 8080 (#0)
> GET /games HTTP/1.1
> Host: localhost:8080
> User-Agent: curl/7.81.0
> Accept: */*
> < HTTP/1.1 405 Method Not Allowed
< Content-Length: 0
<
* Connection #0 to host localhost left intact
\end{shellcode}

这回像样了。以 \texttt{>} 开头的行是客户端发送的头数据；以 \texttt{<} 开头的行是服务返回的头数据。这里首先要注意的是第一条发送的头数据——我们确实发送了 GET 方法。如果没有显式提供动词，\texttt{curl} 发送请求时会使用默认动词——本例中就是 GET。这条输出里另一处有意思的行是响应部分的第一行：可以看到服务器返回了 405 状态码，以及它的明确含义 \texttt{Method Not Allowed}——正是我们所期待的。

\texttt{curl} 允许通过 \texttt{-\-request}（或 \texttt{-X}）选项指定方法。在 Postman 里，你只需通过下拉框改变所用的方法。而在 web 浏览器里，事情就变得有点棘手了。有时可以借助开发者设置发送 POST——或 PUT、DELETE 或任何你想要的方法——但大多数情况下，浏览器能为我们的目的提供的功能已经到头了。从现在起，我们把外部工具测试的范围限定在 \texttt{curl}——Postman 对所有基本用法都相当直观，不需要指南。往我们的端点射一个 POST：

\begin{shellcode}[]
> curl -v -X POST localhost:8080/games
\end{shellcode}

输出如下：

% SOURCE_PDF_PAGE: 073
\begin{shellcode}[]
*   Trying 127.0.0.1:8080...
* Connected to localhost (127.0.0.1) port 8080 (#0)
> POST /games HTTP/1.1
> Host: localhost:8080
> User-Agent: curl/7.81.0
> Accept: */*
>
< HTTP/1.1 200 OK
< Content-Length: 19
< Content-Type: text/plain; charset=utf-8
<
* Connection #0 to host localhost left intact
Creating a new game
\end{shellcode}

这几乎正是我们想要的。说``几乎''，是因为 HTTP 有一条规则：创建资源的动作应当返回一个说明``资源已创建''的状态。我们在这对交换的头里注意到：响应的状态码是 200，代表 OK。创建游戏时不应使用这个状态码；按照 HTTP 标准，我们应当使用 201，代表 Created。表 8.1 描述了来自 HTTP Semantics RFC9110 文档（\url{https://mng.bz/yWPG}）中最常见的 HTTP 状态码。这份文档包含关于 HTTP 协议、标准和状态码的全部细节。记得去查查 418 的含义以及它的由来。

\Needspace{0.46\textheight}
\medskip
\noindent\textbf{表 8.1\quad 最常见的 HTTP 状态码}

\begin{center}
\small
\begin{tabularx}{\textwidth}{|l|X|}
\hline
\textbf{状态码} & \textbf{含义} \\
\hline
\texttt{200} & OK——服务器返回客户端所请求的内容。 \\
\hline
\texttt{201} & Created——服务器处理了请求，并把新创建的资源作为响应发回。 \\
\hline
\texttt{202} & Accepted——服务器将在之后处理客户端的请求——通常用于异步流程。 \\
\hline
\texttt{204} & OK，没有可返回的内容——服务器正确完成了请求，但没有内容可返回。它被 POST、PUT 和 DELETE 命令使用。 \\
\hline
\texttt{400} & Bad Request——服务器因认为存在客户端错误（例如请求格式错误、缺少必填字段）而无法处理。 \\
\hline
\texttt{401} & Unauthorised——客户端不被允许执行该操作；请求缺少目标资源的有效身份验证凭据。 \\
\hline
\texttt{403} & Forbidden——客户端需要身份验证才能访问该资源；换句话说，服务器理解了这个请求，但拒绝满足它。 \\
\hline
\texttt{404} & Not Found——服务器没有找到目标资源的当前表示，或者不愿透露它存在。 \\
\hline
\texttt{500} & Internal Error——服务器遇到了错误，不知道如何处理。 \\
\hline
\end{tabularx}
\end{center}

% SOURCE_PDF_PAGE: 074
在收尾并前往下一个端点之前，把这最后一处修改带进代码吧。负责写出这个状态码的是处理器。那么，打开 \texttt{newgame/handler.go}，加入写出状态码的调用，如代码清单 8.8 所示。这段代码出现在响应中，我们使用 \texttt{WriteHeader} 方法——它只接受一个参数，即随响应体一起携带的状态码——文件是 \texttt{internal/handlers/newgame/handler.go}。

% SOURCE_PDF_PAGE: 075
\begin{gocode}[title={代码清单 8.8 \texttt{handler.go}：设置成功响应的状态码}]
func Handle(w http.ResponseWriter, req *http.Request) {
    [...]
    _, _ = w.Write([]byte("Creating a new game"))
    w.WriteHeader(http.StatusCreated)
}
\end{gocode}

现在，如果重启服务器并运行

\begin{shellcode}[]
> curl -v -X POST http://localhost:8080/games
\end{shellcode}

我们应该会看到响应状态码变成 201，不是吗？唉，很遗憾，并不是。还是 200。这是怎么回事？好吧，看看运行服务器的那个终端，我们应该能看到类似这样的一行：

\begin{shellcode}[]
http: superfluous response.WriteHeader call from
learngo-pockets/httpgordle/internal.newGameHandler (newgamehandler.go:16)
\end{shellcode}

这是什么意思？我们只调用了一次 \texttt{WriteHeader}，怎么会``多余''？原来，\texttt{w.Write} 调用已经先行设置了一个带默认 \texttt{http.StatusOK} 状态码的头。好消息是，一旦头被设置，\texttt{WriteHeader} 会忽略之后的所有调用，所以这基本上就是一个普通的顺序 bug——\texttt{w.Write} 与 \texttt{w.WriteHeader} 谁先调用，谁决定状态码。这是一场被操纵的比赛，因为负责的是我们，由我们控制谁先被调用。让我们调整前面的代码来结束本节，文件仍是 \texttt{internal/handlers/newgame/handler.go}。

% SOURCE_PDF_PAGE: 076
\begin{gocode}[title={代码清单 8.9 \texttt{handler.go}：设置响应消息}]
func Handle(w http.ResponseWriter, req *http.Request) {
    if req.Method != http.MethodPost {
        w.WriteHeader(http.StatusMethodNotAllowed)
        return
    }
    w.WriteHeader(http.StatusCreated)
    _, _ = w.Write([]byte("Creating a new game"))
}
\end{gocode}

太好了！重启之后，我们可以观察到服务器现在的行为符合预期：它返回正确的状态码，并告诉我们它正在创建一局游戏。我们已经验证它只接受 POST 请求。那行关于 \texttt{WriteHeader} 多余调用恼人的提示也随之消失了。我们可以愉快地提交并继续前进，不过，有一种办法能让代码更短。

\subsubsection*{开源多路复用器}

在上一节里，我们做了一个假设：假设端点列表一成不变，永远不需要再实现别的东西。这意味着我们可以为每个端点各配一条路径。但现在假设我们得到通知：需要一个新端点，用来跟踪正在进行的游戏。最合适的 URL 是 \texttt{/games}，方法是 GET。在 Go 1.22 之前，麻烦就出在这儿。过去 Go 每条路径只接受一个处理器，不动词一概不论，这意味着我们不能把同一个 \texttt{/games} 同时用于创建新游戏和获取玩家会话。

所幸，我们现在可以用标准库实现需求。正因如此（以及一些我们会在后文遇到的其他原因），自己动手实现 \texttt{http.Server} 接口一度相当流行，GitHub 上星标最多的 Go 项目里有不少都是干这个的。下面是一些流行的开源库，它们实现的多路复用器支持比 \texttt{net/http} 包曾经提供的更多：

% SOURCE_PDF_PAGE: 077
\begin{description}
  \item[\texttt{github.com/go-chi/chi}] 这个库让端点的实现非常简单。
  \item[\texttt{github.com/gorilla/mux}] 这个库非常完整、健壮。
  \item[\texttt{github.com/gin-gonic/gin}] 这个库是 GitHub 上最受欢迎的 mux。
\end{description}

比如说，如果你想用 chi，需要执行以下步骤。首先，获取依赖。为此，用下面的命令告诉 go 工具把它加进 \texttt{go.mod} 文件：

\begin{shellcode}[]
> go get -u "github.com/go-chi/chi/v5"
\end{shellcode}

你可能注意到这里有个 \texttt{v5} 后缀（用他们仓库里的最新版本即可）。如果你试图在浏览器里访问这个 URL，会返回错误。然而，这个 GitHub 仓库在某个时间点被打上了 \texttt{v5.0.0} 标签（并且随时间推移会有更多标签），所以在 \texttt{go get} 命令里使用 \texttt{/v5} 能确保版本与 v5 兼容——可能是 \texttt{v5.0.0} 或 \texttt{v5.0.8}；两者提供相同的 API。

这些库各有千秋。有的支持成千上万的并行请求，有的提供极快的响应。但我们认为，既然 Go 标准库已经提供了当初促成这些库诞生的那个特性，那就不妨讲讲标准库的用法。这些开源库虽然流行，却无法保证你读到这几行字时它们还在维护。事实上，就在我们写作本书时，\texttt{gorilla/mux} 的维护一度暂停、随后又恢复了。那次支持中断正是促使 Go 团队在标准库中实现更完整版本的诱因之一。

% SOURCE_PDF_PAGE: 078
标准库之中。于是我们可以更新路由器，把方法直接写进 \texttt{HandleFunc} 调用里。

\begin{gocode}[title={代码清单 8.10 \texttt{internal/handlers/router.go}：设置动词与路径}]
package handlers

import (
    "net/http"
    "learngo-pockets/httpgordle/internal/api"
    "learngo-pockets/httpgordle/internal/handlers/newgame"
)

// NewRouter returns a router that listens for requests
// to the following endpoints:
//   - Create a new game;
//
// The provided router is ready to serve.
func NewRouter() *http.ServeMux {
    r := http.NewServeMux()
    r.HandleFunc(http.MethodPost+" "+api.NewGameRoute, #1
        newgame.Handle)  #1
    return r
}

#1 Registers the NewGame endpoint with the POST verb
\end{gocode}

\begin{codeannotations}
\item 以 POST 动词注册 \texttt{NewGame} 端点
\end{codeannotations}

现在路径已经成为端点模式的一部分，我们可以在 NewGame 处理器中移除方法检查，如下一个代码清单所示。我们甚至可以顺势真正返回一个游戏——按 API 的定义来。

\begin{gocode}[title={代码清单 8.11 \texttt{newgame/handler.go}：返回一个 Game 响应}]
func Handle(w http.ResponseWriter, req *http.Request) {
    w.Header().Set("Content-Type", "application/json") #1
    w.WriteHeader(http.StatusCreated)
    apiGame := api.GameResponse{} #2
    err := json.NewEncoder(w).Encode(apiGame) #3
    if err != nil { #4
        // The header has already been set. Nothing much we can do here.
        log.Printf("failed to write response: %s", err)
    }
}

#1 Tells the consumer that we're sending JSON
#2 Defines an empty response
#3 Encodes the game into JSON
#4 If the encoding failed, it's an internal error.
\end{gocode}

\begin{codeannotations}
\item 告知使用方我们在发送 JSON
\item 定义一个空响应
\item 把游戏编码为 JSON
\item 如果编码失败，那是内部错误
\end{codeannotations}

% SOURCE_PDF_PAGE: 079
试一试吧。现在需要测试这个处理器。

\subsubsection*{测试游戏创建}

这个更短的处理器新版本应当更容易测试：代码行数越少，bug 越少。可能的障碍是它需要两个相当复杂的参数，我们得对它们进行 mock 或 stub。

写服务的人很多，所以人们开发了好几个库来让实现服务更容易；测试这些库的人也很多，所以我们不需要亲自去 stub 它们。把测试限制在你自己写的代码上，永远是件好事。

用 \texttt{http.NewRequest} 可以轻松创建请求。至于写入器，Go 有内置的 \texttt{httptest} 包，带一个 Recorder 类型和 \texttt{NewRecorder} 构造函数。

我们要用上两个非常常用的测试包：\texttt{require} 和 \texttt{assert}。使用它们可能引发争论，因为一些纯粹主义者会主张只调用标准库；尽管如此，我们还是用它们，好让你熟悉它们。它们位于模块 \texttt{github.com/stretchr/testify}，先用 \texttt{go get github.com/stretchr/testify} 把这个模块加进项目。

\begin{infobox}{testify 库：assert 与 require}
testify 库在 Go 代码测试中非常流行。它提供了大量校验工具，最常用的是它的两个包 \texttt{assert} 和 \texttt{require}。两个包提供的函数相似——检查错误是否为 nil、两个值是否相等——但行为差别很大：
% SOURCE_PDF_PAGE: 080
\begin{itemize}
  \item 如果 \texttt{assert} 包中的函数发现异常，会显示一条消息，测试继续执行。
  \item 如果 \texttt{require} 包中的函数发现异常，会显示一条消息，测试立即终止。
\end{itemize}

这帮助我们决定想用哪个包：需要检查多个值时用前者；明知出错后继续测试已无意义时用后者。现在你手头应当一应俱全，可以为正常行为（nominal behavior）编写测试了。
\end{infobox}

\begin{gocode}[title={代码清单 8.12 \texttt{newgame/handler\_test.go}：测试 \texttt{Handle}}]
func TestHandle(t *testing.T) {
    req, err := http.NewRequest(
        http.MethodPost, "/games", nil) #1
    require.NoError(t, err)
    recorder := httptest.NewRecorder() #2
    Handle(recorder, req) #3
    assert.Equal(t, http.StatusCreated, recorder.Code) #4
    assert.Equal(t, "application/json",
        recorder.Header().Get("Content-Type")) #4
    assert.JSONEq(t, `{"id":"",...}`,
        recorder.Body.String()) #4
}

#1 Creates a request
#2 Creates a response recorder
#3 Calls the function
#4 Asserts that the response is as
expected
\end{gocode}

\begin{codeannotations}
\item 创建请求
\item 创建响应记录器
\item 调用该函数
\item 断言响应符合预期
\end{codeannotations}

这是我们的第一个端点——我们的服务现在支持创建（空的）游戏了。这是重大一步，我们涵盖了 web 服务的许多重要方面。在开始游戏之前，下一个任务是确保能够根据标识符获取一局游戏的状态。

% SOURCE_PDF_PAGE: 081
\subsection{获取游戏状态}

到目前为止的情况是：我们拥有一个可以创建 Gordle 游戏的服务。当然，最终目标是让玩家进行猜测，但这里要描述的第二个端点是 \texttt{GetStatus}。它只比第一个多一点东西，因为只引入了一个新概念：从用户那里读取可变的输入。这一次，我们不能总是返回``一局游戏''了——我们需要能识别用户想看哪一局。

\subsubsection*{向 HTTP API 传参}

用户向 HTTP web 服务传递参数（或变量）主要有四种方式。我们会看到，它们各自适用于特定的用例：

\begin{itemize}
  \item 路径参数（path parameter）
  \item 查询参数（query parameter）
  \item 请求体（request body）
  \item 元参数（metaparameter）
\end{itemize}

当我们想定位单个资源时，使用路径参数。在我们的 Gordle 服务器里，可以用一个标识符来定位一局游戏的实例。定位一局游戏的路径应当是 \texttt{/games/\{gameID\}}。这里的曲线花括号意味着我们给路径参数赋予了一个通配标识符——不同用户调用我们的端点时 \texttt{\{gameID\}} 的值各不相同。

在 REST API 中使用 \texttt{/items/\{itemID\}} 是常见实践（比单数版本 \texttt{/item/\{itemID\}} 更常见）。有时我们可以接受多于一个路径参数。比如 Twitter 用 \texttt{https://twitter.com/\{user\}/status/\{messageID\}} 展示某位用户的一条消息。类似地，Wikipedia 用路径

% SOURCE_PDF_PAGE: 082
参数来访问它的文章：\texttt{http://jp.wikipedia.org/wiki/金継ぎ}，其中文章的标识符就是金継ぎ。我们的 \texttt{GetStatus} 端点将为 Gordle 服务实现同样的机制。

查询参数用于筛选请求想要定位的资源。筛选器是键值对的列表。结果可能有零条、一条或许多条——我们不知道，也不能做任何假设。这些查询参数放在请求的 URI 中，但位于末尾，用一个 \texttt{?} 字符与 URL 分隔。这些参数并不专属于某个端点，可以在 API 的多个地方使用。其语法是 \texttt{\{path\}?key1=value1\&key2=value2}。最常见的例子是 Google 的搜索引擎：人们向 Google 提出的每个可能的查询，不可能（合理地）都有专门的资源，所以每个查询都作为查询参数发送到他们的服务器，键是 \texttt{q}，值由用户键入：\texttt{www.google.com/search?q=金継ぎ}。本章末尾我们会演示如何用查询参数改进 \texttt{NewGame} 端点，让调用者指定想用的语言。

路径参数和查询参数应当用于指定我们想操作（获取、删除、更新）哪些资源，或这些资源应具备什么特征；而有时我们需要提供请求本身固有的参数。当我们需要向服务发送数据时，使用体参数（body parameter）。这个名字来源于它们会作为请求体的一部分传输给服务。到目前为止我们还没见过请求体，而这正是第三个端点——Guess——的重点。

最后，有些参数是元参数——它们不影响执行本身，但会影响
% SOURCE_PDF_PAGE: 083
输出的格式，或描述调用者的一些信息。这些参数放在查询的请求头中——正如状态码放在响应的头中一样。请求头通常是我们存放身份验证信息的地方。理论就讲到这里，开始实现新端点吧！

\subsubsection*{为 GETSTATUS 定义 HTTP API}

与前面一样，我们需要声明这个新端点的路径，以及调用时期望的方法。我们把这两个值放在 \texttt{api} 包里，让其他使用者可见。我们想获取游戏资源的状态而不改变它，所以 GET 就够了。但路径稍微复杂一点！在 REST API 中，每个请求都必须包含识别其目标资源的全部信息。在我们的例子里，这意味着请求需要包含游戏的 ID。如前所述，路径参数是提供这个标识符的常见方式：\texttt{/games/\{game\_id\}}。

如何表示一条不恒定的路径？这正是默认 \texttt{net/http} 包过去过于严苛之处，也是促使开发者编写自己的路由库（如 chi）的问题之一。我们希望能够通过路径 \texttt{/games/8476516} 访问游戏状态，其中 \texttt{8476516} 是游戏的标识符。显然，我们不可能创建数十亿条路由——每个标识符一条——所以，我们改为让标准库利用路由中的曲线花括号来决定路径参数如何处理。

在开始监听 \texttt{GetStatus} 路由之前，还有最后一个定义要写进文档：期望的响应状态码是什么？这个请求是在
% SOURCE_PDF_PAGE: 084
索取一个资源，所以可能的响应应当是``200，拿去吧''或——如果游戏不存在——``404，没找到''。当然，``500，服务器内部错误''永远有可能发生，但我们要尽力避免它。

\subsubsection*{实现 GETSTATUS}

为新处理器创建一个包，理想情况下是 \texttt{internal/handlers/getstatus}，其中为主要的 \texttt{Handle} 函数建一个文件。如果你认定从 \texttt{newgame} 包复制粘贴是个好选择，别忘了在两个新文件里重命名包名（如果你尽职地添加了 \texttt{doc.go}，那就是三个文件）。

目前，\texttt{GetStatus} 端点的实现将仅限于打印调用者作为路径参数传入的游戏标识符：我们还没有任何存储，只是想确保自己知道如何解析 ID。先更新 \texttt{api} 包，加入新端点的路径。

\begin{gocode}[title={代码清单 8.13 \texttt{api/http.go}：添加新端点路由}]
const (
    // GameID is the name of the field that stores the game's identifier
    GameID = "id"
    // NewGameRoute is the path to create a new game.
    NewGameRoute = "/games"
    // GetStatusRoute is the path to get the status
    // of a game identified by its id.
    GetStatusRoute = "/games/{" + GameID + "}"
}
\end{gocode}

在这个文件里首先要注意的是：我们为游戏 ID 定义了一个常量。我们不只是想定义 \texttt{GetStatus} 端点的路由——以后还想取回那个路径参数，而这需要用到 \texttt{GameID} 常量。

% SOURCE_PDF_PAGE: 085
接下来，把路径添加到路由器。向 mux 添加处理器没有特定的优先级。按资源和相关行为分组通常是最合理的：

\begin{gocode}[]
...
    r.HandleFunc(http.MethodPost+" "+api.NewGameRoute,
        newgame.Handle) #1
    r.HandleFunc(http.MethodGet+" "+api.GetStatusRoute,
        getstatus.Handle) #2
...

#1 Previous endpoint
#2 Adds the new endpoint
\end{gocode}

\begin{codeannotations}
\item 之前的端点
\item 添加新端点
\end{codeannotations}

现在需要编写这个 \texttt{getstatus.Handle} 函数。它必须与第一个端点有相同的签名，因为它是一个 \texttt{HandlerFunc}。要取回路径参数，可以使用 \texttt{request.PathValue(api.GameID)}。如果出了任何问题，我们会调用 \texttt{http.Error}，它会向响应写入器写一条消息和一个状态码。记住：调用 \texttt{http.Error} 之后总要 return，以防止再向响应写入任何内容。让我们写到 \texttt{internal/handlers/getstatus/handle.go} 文件中。

\begin{gocode}[title={代码清单 8.14 \texttt{handle.go}：状态端点处理器}]
func Handle(
    w http.ResponseWriter, request *http.Request) { #1
    id := request.PathValue(api.GameID) #2
    if id == "" { #3
        http.Error(writer, "missing the id of the game",
            http.StatusBadRequest)
        return
    }
    log.Printf("retrieve status of game with id: %v", id)
    apiGame := api.GameResponse{
        ID: id,
    }
    // ... encode into JSON
}

#1 Follows the required signature
#2 Parses the parameter
#3 Checks that the ID isn't empty
\end{gocode}

\begin{codeannotations}
\item 遵循要求的签名
\item 解析参数
\item 检查 ID 不为空
\end{codeannotations}

为了简单起见，我们决定使用标准的 \texttt{log} 包，它不是线程安全的。所以记住，

% SOURCE_PDF_PAGE: 086
这可能导致日志乱序，并给后续测试添堵。ID 的校验本可以更彻底：我们可以检查它只含数字，或符合我们出于安全考虑设下的任何约束。为了让项目保持``袖珍尺寸''，我们先不管这个；但如果你要把东西推上生产环境，请把它记在心里。

运行服务器，调用这个端点。它返回你在 URL 里传入的 ID 了吗？恭喜——可以提交了。等等，那测试呢？幸好你问了。

\subsubsection*{测试 GETSTATUS}

这里的测试与 newGame 版本差别不大。我们只需要把 ID 加进 URL 参数列表。要创建一个带路径参数的请求，先创建一个指向我们端点的请求——使用正确的动词和路由——然后设置路径值。

\begin{gocode}[title={代码清单 8.15 \texttt{getstatus/handler\_test.go}：测试 \texttt{handle}}]
func TestHandle(t *testing.T) {
    req, err := http.NewRequest(http.MethodGet,
        "/games/", nil) #1
    require.NoError(t, err)
    // add path parameter
    req.SetPathValue(api.GameID, "123456")
    recorder := httptest.NewRecorder() #2
    Handle(recorder, req)
    assert.Equal(t, http.StatusOK, recorder.Code)
    assert.JSONEq(t, `{"id":"123456","attempts_left":0,"guesses":[],
        "word_length":0,"status":""}`, recorder.Body.String())
}

#1 Using "/games/ 1 23456" wouldn't work here.
#2 Creates the response receiver
\end{gocode}

\begin{codeannotations}
\item 这里用 ``/games/ 1 23456'' 行不通
\item 创建响应接收器
\end{codeannotations}

其余的你可以自己猜。现在我们既能创建游戏，也能取回它们。这让我们确信，万事俱备，可以迎接第三个也是最后一个端点——猜词！

% SOURCE_PDF_PAGE: 087
\subsection{猜词}

最终，要玩游戏，玩家必须能带着自己的单词发起查询，并得到一条反馈消息。第三个端点的 API 会是什么样？

\subsubsection*{请求定义}

添加新端点意味着选择一对新的``路径 + 方法''。因为我们在修改一个已存在的资源，所以 PUT 合适。路径相当直白——\texttt{/games/\{game\_id\}}，与 getStatus 端点类似，因为这就是我们要交互的资源。

但这回端点需要接收参数。它将从请求体中读取猜测；由于我们遵循 HTTP 和 REST 标准，编码为 JSON。某些情况下，更新一个资源需要发送它的完整描述。那样的话，POST 和 GET 的响应、PUT 的请求就会共用同一个 JSON 结构。这里，修改一局游戏的状态只需要发送一个单词——只要提供了 ID——所以我们用最简单的 JSON 对象：

\begin{gocode}[]
{"guess":"hello"}
\end{gocode}

把它翻译成一个 Request，定义在 \texttt{api} 包里供他人使用。

\begin{gocode}[title={代码清单 8.16 \texttt{api/guess.go}：请求定义}]
// GuessRequest is the structure of the message used when submitting a guess.
type GuessRequest struct {
    Guess string `json:"guess"` #1
}

#1 Uses JSON tags for encoding
\end{gocode}

\begin{codeannotations}
\item 使用 JSON 标签进行编码
\end{codeannotations}

% SOURCE_PDF_PAGE: 088
把端点的路径也加进这个文件。路径与 GetStatus 相同，但没有任何东西保证永远如此，所以我们需要另一个常量。共用同一个路径常量，等于从设计上强迫二者永远一致。我们不想那样，因为每个端点都应当拥有自己的路径常量：

\begin{gocode}[]
GuessRoute = "/games/{" + GameID + "}"
\end{gocode}

我们真正想从设计上强制的，是 \texttt{GameResponse} 结构的一致性。GetStatus 和 Guess 端点都返回完整的游戏，我们不希望两者的定义渐行渐远。为此，两个端点直接使用同一个结构作为响应即可。

\subsubsection*{解码请求体}

API 已经有了，是时候在新包里写一个新处理器并接入路由器了。这里没有需要新解释的东西，放手准备 handle 函数吧。你甚至可以先加一行简单的 \texttt{log.Printf} 跑一跑，确保它如预期工作。

这个新处理器首先解析游戏的 ID，与上一个端点完全一样，不足为奇。其次，它要解析请求体。你在前几章已经知道如何把 JSON 消息解码到 Go 结构体。某位天才灵光一闪：\texttt{http.Request} 的 \texttt{Body} 字段应当实现 \texttt{io.Reader} 接口——那我们就好好利用这一点。

% SOURCE_PDF_PAGE: 089
\begin{gocode}[title={代码清单 8.17 \texttt{guess/handler.go}：如何读取请求体}]
// Read the request, containing the guess, from the body of the input.
r := api.GuessRequest{} #1
err := json.NewDecoder(request.Body).Decode(&r) #2
if err != nil {
    http.Error(writer, err.Error(),
        http.StatusBadRequest) #3
    return
}

#1 Instantiates the request struct
#2 Decodes the body
#3 If the JSON doesn't parse, it's a bad request.
\end{gocode}

\begin{codeannotations}
\item 实例化请求结构体
\item 解码请求体
\item 如果 JSON 无法解析，就是一个错误请求
\end{codeannotations}

把结果打印到日志里，然后用 \texttt{curl} 射一下服务，看看会发生什么（使用 \texttt{curl} 时，路径参数不用特殊指定，直接写就行）：

\begin{shellcode}[]
> curl -v -X PUT "http://localhost:8080/games/123456" -d '{"guess":"hello"}'
\end{shellcode}

注意 \texttt{-d} 标志用于传递请求体。如果你的请求体保存在文件里，也可以用 \texttt{-d@file.json}。日志里的一切显示正常吗？测试写了吗？

\subsubsection*{用请求体做测试}

如果我们对这个 handle 函数运行与 GetStatus 相同的测试，它会 panic：我们的请求体是 nil，等于在从一个 nil reader 解码，结局不会好。需要做的唯一修改相当短：我们一直用来创建请求的 \texttt{http.NewRequest} 函数，其第三个参数就是一个 \texttt{io.Reader}，在 Go 里从字符串创建一个 reader 很容易（记住，用反引号 \texttt{`} 包裹含有双引号 \texttt{"} 的字符串，就不必转义）：

\begin{gocode}[]
body := strings.NewReader(`{"guess":"pocket"}`)
req, err := http.NewRequest(http.MethodPut, "/games/", body))
req.SetPathValue(api.GameID, "123456")
\end{gocode}

% SOURCE_PDF_PAGE: 090
现在我们拥有了服务的完整结构：三个端点，各能返回一些东西。这是把服务部署到测试环境、让别人把玩一番的好时机。此时的典型画面是：你把焕然一新的服务链接发给团队其他人，附上一大段话请他们试用，提醒他们它还什么都不会干，然后等某人回复说``嘿，干得漂亮，不过我发现一个 bug：为什么它总是返回一个空游戏？''好，来修掉它。

\section{领域对象}

每个玩家的游戏必须存在某个地方，通常称之为仓库（repository，简称 repo）。在众多选项中，最便宜也最快的是存在内存里。这样做有很多缺点：如果你决定稍微扩展一下、部署多个实例，游戏只存在其中一个实例上，这意味着如果你的 Guess 查询打到一台没有注册这局游戏的实例上，你会得到 404。此外，如果那个实例因任何原因宕机，你的游戏就丢了。

我们想让这个项目能装进口袋，所以暂时选这条路。在更大的项目里，我们会用一个正经的数据库。

\subsubsection*{关注点分离}

判断一个项目的代码是否整洁，主要方法之一就是看关注点是否分离。理论上，每个包、每个结构体都应当有一个能用一句话说明的角色，而这句话就是其文档的第一行。大多数时候，一句话足以概括全部。实践中，这条规则可能在高度
% SOURCE_PDF_PAGE: 091
耦合的想法之间引入复杂的通信。如前所述，在``拆开''与``放在一起''之间总有权衡要找。

如果一个包的职责无法用寥寥数语概括，这通常是个危险信号：未来的维护者将不知道去哪里找哪段逻辑，最后干脆把一切扔掉重写——而且不一定重写得更好。有时，我们重写代码只是为了搞懂到底发生了什么。

无论你选择六边形架构、千层面式 MVC，还是任何其他适合你需求的奇美拉，大多数时候，你都想把数据存储管理与 API 细节隔离开来。你确实会见到一些数据存储与 API 设计需要协同演化才能保证性能的案例，但那是少数。

软件理论家们常常举这样一个完美数据存储包的例子：``看，你在用 MariaDB，只需要改一下这个包的导入，砰，你就在用 DynamoDB 了。''想法很美好，但它从未发生：没有人会简简单单地更换存储系统。人们经常做的是维护它，而``所有与数据库有关的东西都在这里、所有与数据库无关的东西都不在这里''这一保证，长远来看对每个人都有帮助。

\subsection{领域类型}

你可以为数据仓库新建一个包。你想让其他模块——其他开发者——使用它吗？不想。那 \texttt{internal/repository} 就够了。这个包负责存储和取回游戏。就这里而言，这就是一句话的文档。图 8.1 展示了到目前为止代码的结构。

% SOURCE_PDF_PAGE: 092
\begin{center}
\begin{tikzpicture}[node distance=15mm,
  box/.style={draw=PocketTealDark, thick, rounded corners=2.2mm, fill=PocketMist,
              minimum width=30mm, minimum height=14mm, font=\ttfamily\small, align=center}]
  \node[box] (api) {API};
  \node[box, below=of api] (handlers) {Handlers};
  \draw[-{Stealth[length=2.6mm]}, thick] (handlers) -- (api);
\end{tikzpicture}

\medskip
\noindent\textbf{图 8.1 服务的结构}
\end{center}

我们有一个可被其他模块导入的 \texttt{api} 包，以及一个由 main 直接调用的内部 \texttt{handlers} 包。若要以六边形架构（又称端口与适配器）的心态来进行这一开发，我们应当朝着图 8.2 所示的设计努力：数据访问仓库、Gordle 库与处理器各自独立工作，以领域（domain）作为共同语言。

% SOURCE_PDF_PAGE: 093
\begin{center}
\begin{tikzpicture}
  % central hexagon (adapters ring)
  \node[regular polygon, regular polygon sides=6, minimum width=64mm, minimum height=48mm,
        draw=PocketInk, fill=PocketGreen!14, thick] (hex) at (0,0) {};
  \node[regular polygon, regular polygon sides=6, minimum width=26mm, minimum height=20mm,
        draw=PocketInk, fill=PocketGreen!70!black, thick, font=\sffamily\bfseries\small\color{white}] at (0,0) {Domain};
  \node[font=\footnotesize, fill=PocketGreen!14, inner sep=1pt] at (-2.30,0) {Adapters};
  \node[font=\footnotesize, fill=PocketGreen!14, inner sep=1pt] at (2.30,0) {Adapters};
  \node[font=\footnotesize, fill=PocketGreen!14, inner sep=1pt] at (0,-1.90) {Adapters};
  % outer boxes
  \node[draw=PocketTealDark, thick, rounded corners=1.8mm, fill=PocketTeal!22,
        minimum width=25mm, minimum height=11mm, font=\ttfamily\small] (api) at (-5.3,1.75) {API};
  \node[draw=PocketTealDark, thick, rounded corners=1.8mm, fill=PocketTeal!22,
        minimum width=25mm, minimum height=11mm, font=\ttfamily\small] (handlers) at (-5.3,0) {Handlers};
  \node[draw=PocketTealDark, thick, rounded corners=1.8mm, fill=PocketTeal!22,
        minimum width=27mm, minimum height=11mm, font=\ttfamily\small] (repo) at (5.4,0) {Repository};
  \node[draw=PocketTealDark, thick, rounded corners=1.8mm, fill=PocketTeal!22,
        minimum width=29mm, minimum height=11mm, font=\ttfamily\small] (gordle) at (0.1,-3.30) {Gordle library};
  % arrows
  \draw[-{Stealth[length=2.6mm]}, thick] (handlers) -- (api);
  \draw[-{Stealth[length=2.6mm]}, thick] (hex.west) -- (handlers);
  \draw[-{Stealth[length=2.6mm]}, thick] (hex.east) -- (repo);
  \draw[-{Stealth[length=2.6mm]}, thick] (0,-2.42) -- (gordle);
\end{tikzpicture}

\medskip
\noindent\textbf{图 8.2 服务的目标结构}
\end{center}

对于这个服务的尺寸来说，这有点过于复杂，所以我们把领域逻辑留在 \texttt{handlers} 包里。不过要记住：核心逻辑与 API 细节应当是两件不同的事，中间有一条容易画出的边界。我们设计中的每一层都应为特定目的服务。

我们这里所说的领域（domain），是服务的核心，即程序的业务逻辑。所有适配器都应当能够依赖领域，而领域不依赖它们中的任何一个。这样，我们就能避免代码中的循环依赖，也避免大脑里的死结。

这个包可以简单叫作 \texttt{domain}、\texttt{core}，或者取一个更具体但受限的名字。就我们的情况而言，我们发现它处理的是玩家的会话。

在 \texttt{internal/session} 包中，创建一个 \texttt{Game} 结构体。它应当包含服务与 Gordle 游戏交互所需的一切。我们知道，我们希望用一个 \texttt{Game} 结构来响应 API 的需求。

% SOURCE_PDF_PAGE: 094
\begin{gocode}[title={代码清单 8.18 \texttt{internal/session/game.go}：定义 \texttt{Game} 结构}]
// Game contains the information about a game.
type Game struct {
    ID           GameID
    AttemptsLeft byte
    Guesses      []Guess
    Status       Status
}
\end{gocode}

如清单所示，我们需要一些新类型，那就把 \texttt{GameID} 和 \texttt{Status} 定义在刚才同一个文件里。

\begin{gocode}[title={代码清单 8.19 \texttt{game.go}：定义标识符与状态类型}]
// A GameID represents the ID of a game.
type GameID string #1

// Status is the current status of the game and
// tells what operations can be made on it.
type Status string #B

const (
    StatusPlaying = "Playing" #2
    StatusWon     = "Won"
    StatusLost    = "Lost"
)

#1 Game ID is a simple string.
#2 Enum-like type and its list of constants
\end{gocode}

\begin{codeannotations}
\item 游戏 ID 就是一个简单的字符串
\item 枚举风格的类型及其常量列表
\end{codeannotations}

我们还希望暴露猜测，因为需要带着它们到处走、存储它们、返回它们。下面的代码清单可以做到。

\begin{gocode}[title={代码清单 8.20 \texttt{game.go}：定义 \texttt{Guess} 类型}]
// A Guess is a pair of a word (submitted by the player) and
// its feedback (provided by Gordle).
type Guess struct {
    Word     string
    Feedback string
}
\end{gocode}

这是个好的开始。还不足以玩游戏，但那是下一节的事。

有一件事已经可以加进领域里：一个业务错误。如果玩家在游戏
% SOURCE_PDF_PAGE: 095
结束后又发来一次猜测，我们应当明确指出问题。你可以定义一个自定义错误类型，或者走捷径：

\begin{gocode}[]
// ErrGameOver is returned when a play is made but the game is over.
var ErrGameOver = errors.New("game over")
\end{gocode}

本书的一位作者非常不喜欢暴露全局变量，但你自己决定就好。

\subsection{API 适配器}

领域实体已经定义好了，我们可以把它接入六边形架构的各个部分。从 API 这一侧开始。

\subsubsection*{NEWGAME}

以第一个端点 \texttt{NewGame} 为例。处理器目前是这样创建游戏 API 版本的：

\begin{gocode}[]
func Handle(w http.ResponseWriter, req *http.Request) {
    w.Header().Set("Content-Type", "application/json") #1
    w.WriteHeader(http.StatusCreated)  #1
    apiGame := api.GameResponse{} #2
    err := json.NewEncoder(w).Encode(apiGame) #3
    if err != nil //...
}

#1 Encodes the response headers
#2 Creates the game
#3 Encodes the response body
\end{gocode}

\begin{codeannotations}
\item 编码响应头
\item 创建游戏
\item 编码响应体
\end{codeannotations}

\texttt{Handle} 函数的职责是处理 API 细节。如果客户端需要不同的标志或不同的格式，就应当在这里发生。

我们决定把业务逻辑留在同一个包里，但这不意味着要留在同一个函数里。我们想要一个函数负责创建游戏并把它保存
% SOURCE_PDF_PAGE: 096
到某处，另一个函数负责把它改造成客户端想要的形状。代码如下。

\begin{gocode}[title={代码清单 8.21 \texttt{newgame/handler.go}：使用 \texttt{session.Game}}]
func Handle(w http.ResponseWriter, req *http.Request) {
    game, err := createGame() #1
    if err != nil {
        log.Printf("unable to create a new game: %s",
            err) #2
        http.Error(w, "failed to create a new game",
            http.StatusInternalServerError)
        return #3
    }
    w.Header().Set("Content-Type", "application/json")
    w.WriteHeader(http.StatusCreated) #4
    apiGame := response(game) #5
    // ...
}

func createGame() (session.Game, error) {
    return session.Game{}, nil
}

func response(game session.Game) api.GameResponse {
    return api.GameResponse{}
}

#1 Creates the game
#2 Logs any error for future debugging
#3 Don't forget to stop there.
#4 Encodes for HTTP headers
#5 Transforms the game into the API format
\end{gocode}

\begin{codeannotations}
\item 创建游戏
\item 记录任何错误，便于日后调试
\item 别忘了就此打住
\item 为 HTTP 头编码
\item 把游戏转换为 API 格式
\end{codeannotations}

这样，每个函数都有一份定义明确的工作，我们可以测试一个个单元，而不是一坨坨大块头。团队内部的工作也更容易并行，多人做出的修改更容易合并，最重要的是——我们读得懂自己读的东西了。

\subsubsection*{GETSTATUS}

继续 \texttt{GetStatus} 端点。响应格式完全相同，这意味着响应函数可以复用。它应该住在 \texttt{newgame} 包里、再暴露给 \texttt{getstatus} 包吗？对下一个要改这部分代码的人来说，这实在没什么道理。

% SOURCE_PDF_PAGE: 097
这里更直接的解法，是把这个端点放在带有该函数的 \texttt{internal/api} 包里，如代码清单 8.22 所示。我们避免了导入环，接棒的人也能明白里面是什么。

现在 \texttt{response} 这个名字不再合适了。它把一个 \texttt{session.Game} 转换成 \texttt{GameResponse}。它的逻辑大部分也可以先写出来，甚至可以开始测试。

\begin{gocode}[title={代码清单 8.22 \texttt{convert.go}：API 与领域之间的适配器}]
// ToGameResponse converts a session.Game into an GameResponse.
func ToGameResponse(g session.Game) GameResponse {
    apiGame := api.GameResponse{
        ID:           string(g.ID),
        AttemptsLeft: g.AttemptsLeft,
        Guesses:      make([]Guess, len(g.Guesses)), #1
        Status:       string(g.Status),
        // TODO WordLength #2
    }
    for index := 0; index < len(g.Guesses); index++ {
        apiGame.Guesses[index].Word =
            g.Guesses[index].Word #3
        apiGame.Guesses[index].Feedback = g.Guesses[index].Feedback
    }
    if g.AttemptsLeft == 0 { #4
        apiGame.Solution = ""// TODO solution #2
    }
    return apiGame
}

#1 Inits the slice
#2 This is for later; we don't have the data yet.
#3 Accesses the slice elements directly
#4 Shows a solution if the game is over
\end{gocode}

\begin{codeannotations}
\item 初始化切片
\item 这是留给以后的；我们还没有数据
\item 直接访问切片元素
\item 游戏结束时给出答案
\end{codeannotations}

测试它毫无花招：我们需要检查一个输入游戏能以新形状如期输出。可以用 \texttt{go test -cover} 或 IDE 的界面。\texttt{-cover} 选项测量测试执行的代码百分比，帮助你发现代码中任何未经测试的分支或路径。

% SOURCE_PDF_PAGE: 098
。这能确保你没有漏掉边界情况，测试也全面覆盖了程序的功能。

更新 GetStatus 处理器并运行它的测试。还满意吗？如果在上一版里你的猜测切片是 nil、而现在被初始化了，测试就应当失败。你可以通过把 JSON 值 \texttt{null} 替换为空数组 \texttt{[]} 来修复它。

\subsubsection*{GUESS}

终于到达最后一个端点了。看看处理器在这里做什么：

\begin{gocode}[]
func Handle(w http.ResponseWriter, req *http.Request) {
    // ... decode request
    apiGame := api.GameResponse{
        ID: id,
    }
    // ... encode response
}
\end{gocode}

我们已经可以用两次调用来替换这个朴素的 \texttt{api.Game} 初始化：一次调用新的本地（未导出）函数执行业务逻辑，一次转换成 API 结构。所需代码见下一个代码清单。

% SOURCE_PDF_PAGE: 099
\begin{gocode}[title={代码清单 8.23 \texttt{guess/handler.go}：使用 \texttt{session.Game}}]
// Handler returns the handler for the game creation endpoint.
func Handle(w http.ResponseWriter, req *http.Request) {
    id := req.PathValue(api.GameID)
    // ...
    r := api.GuessRequest{}
    // ...
    game, err := guess(id, r) #1
    if err != nil {...}
    apiGame := api.ToGameResponse(game) #2
    w.Header().Set("Content-Type", "application/json")
    // ...
}

func guess(id string, r api.GuessRequest) (session.Game, error) {
    return session.Game{
        ID: session.GameID(id),
    }
}

#1 Plays the game
#2 Adapts for API exposition
\end{gocode}

\begin{codeannotations}
\item 玩这一局游戏
\item 为 API 暴露做适配
\end{codeannotations}

事情的顺序总是一样的：解码请求、校验所有需要校验的输入、调用业务逻辑、把返回的领域类型转换成 API 可读的结构、编码响应。

如果你的测试都通过了，可以提交。试着运行服务、用几发 \texttt{curl} 打一打，看看服务的行为。如果你再次部署到测试环境，那位没读你警告的郁闷队友依然会郁闷，因为游戏还是像以前一样空空如也。

\section{仓库}

此时，我们必须理清优先级。手头还有两大任务：保存游戏，以及让客户端玩游戏。先做哪个？

% SOURCE_PDF_PAGE: 100
做这个决定时，想想哪一个能展示出最多的进展。如果从存储开始，前两个端点就会完工，我们还能在端到端测试中看到它们如何整合（至少是``端到半路''测试，因为我们还玩不了）。如果从玩开始，每次猜测都会在一局新的空游戏上进行，测试会又难又不稳定。看来就该先做存储。

如果我们在使用某种经过验证的数据存储技术，比如 SQL 数据库，就需要一些领域用不上的存储相关额外字段，例如 created at、deleted at、版本控制等等。那样的话，我们会新建一个 \texttt{Game} 结构，并在领域与仓库之间做适配器，就像为 API 做的那样。这样，数据库的 schema 就能独立于领域或 API 自行演化。

既然我们追求的是最快的存储方案，就没有任何技术层面的约束。这意味着可以继续沿用领域结构。

\subsection{内存数据库}

数据库选项多如牛毛，每一种都只在有限场景下理想。如前所述，内存存储在任何场景下都不理想，但它写得快。

在我们的场景里，用一个变量来存游戏。我们对游戏做的所有操作都依赖其 ID，所以键值存储再合适不过。在 Go 中，这就是一个映射。为这个映射建一个包。再问一遍：你想让外部模块使用你的仓库吗？请说不。如果客户端绕过服务直连数据库，那整个服务和它的 API 就失去了存在的意义。你将要添加的任何额外安全或逻辑（发送事件、排行榜等）都会立刻 bug 丛生。

% SOURCE_PDF_PAGE: 101
\subsection{最简单的仓库}

前几章讲过如何创建一个能充当依赖的对象。如果用的是外部数据库，我们会在服务器启动时初始化一条连接，并把它保留为整个服务的依赖。这里，我们改为初始化那个映射，并以同样的方式保留它。

来创建仓库结构。它将持有诸如 \texttt{Find} 和 \texttt{Update} 之类的方法。

\begin{gocode}[title={代码清单 8.24 \texttt{repository/memory.go}：声明结构}]
// GameRepository holds all the current games.
type GameRepository struct {
    storage map[session.GameID]session.Game #1
}

#1 Uses domain types
\end{gocode}

\begin{codeannotations}
\item 使用领域类型
\end{codeannotations}

当然，这需要初始化，所以我们需要一个 \texttt{New()} 函数。代码见下面的清单。

\begin{gocode}[title={代码清单 8.25 \texttt{memory.go}：仓库的 \texttt{New} 函数}]
// New creates an empty game repository.
func New() *GameRepository {
    return &GameRepository{
        storage: make(map[session.GameID]session.Game), #1
    }
}

#1 Inits the map
\end{gocode}

\begin{codeannotations}
\item 初始化映射
\end{codeannotations}

我们本可以有一个同时负责创建和更新的 \texttt{Upsert} 方法。但在我们的场景里，有些特定用例是创建新游戏，另一些只在已有游戏上操作，所以我们倾向于把它们分开。这让我们可以

% SOURCE_PDF_PAGE: 102
验证``我们没有试图把同一个实体创建两次''或``没有往一局已经猜了几次的游戏里插入新游戏''。用下面的代码向存储添加一局游戏。

\begin{gocode}[title={代码清单 8.26 \texttt{memory.go}：向存储添加游戏}]
// Add inserts for the first time a game in memory.
func (gr *GameRepository) Add(game session.Game) error {
    _, ok := gr.storage[game.ID]
    if ok {
        return fmt.Errorf("gameID %s already exists",
            game.ID) #1
    }
    gr.storage[game.ID] = game
    return nil
}

#1 Can't add the same game twice
\end{gocode}

\begin{codeannotations}
\item 不能把同一局游戏加两次
\end{codeannotations}

这个错误我们回头还要说。如果想在调用代码里专门检查它，目前还挺困难。

\begin{infobox}{支线任务}
你可以写出 \texttt{Find} 和 \texttt{Update} 方法。前者需要从映射取回，找不到对应 ID 时返回某种错误。后者同样要防止插入，只允许覆盖已存在的值。

你还可以为这四个函数写单元测试。这里，我们在另外三个测试中直接使用 \texttt{New}，然后视为完工；没有什么特别的窍门。

这个包能用了，也测试过了。可该怎么使用它呢？
\end{infobox}

% SOURCE_PDF_PAGE: 103
\subsection{服务级依赖}

我们希望仓库在启动时初始化，并作为依赖传递给服务路由器。回顾一下我们的 main 现在做什么：

\begin{gocode}[]
func main() {
    err := http.ListenAndServe(":8080", handlers.NewRouter())
    if err != nil {
        panic(err)
    }
}
\end{gocode}

我们可以轻松加入初始化，并把新变量传给路由器。

\begin{gocode}[title={代码清单 8.27 \texttt{main.go}：使用仓库}]
func main() {
    db := repository.New() #1
    err := http.ListenAndServe(
        ":8080", handlers.NewRouter(db)) #2
    if err != nil {
        panic(err)
    }
}

#1 Inits the data storage
#2 Passes the repository to the router
\end{gocode}

\begin{codeannotations}
\item 初始化数据存储
\item 把仓库传给路由器
\end{codeannotations}

路由器如何把它传给处理器？我们的 \texttt{NewRouter} 函数并不调用处理器，它只是把处理器的引用交给路由器，所以没法简单地加个参数。我们能做的，是把 \texttt{Handle} 函数变成启动时创建的匿名函数。

在下面的代码清单里，我们把 \texttt{Handle} 函数匿名化，改用包一层的方式：包在一个 \texttt{Handler} 函数里，它接收一个仓库作为参数，返回之前的 \texttt{http.HandleFunc}。

% SOURCE_PDF_PAGE: 104
\begin{gocode}[title={代码清单 8.28 \texttt{newgame/handler.go}：使用仓库}]
// Handler returns the handler for the game creation endpoint.
func Handler(db *repository.GameRepository) http.HandlerFunc {
    return func(
        w http.ResponseWriter, _ *http.Request) { #1
        game, err := createGame(db)
        // ...
    }
}

#1 Returns a HandlerFunc
\end{gocode}

\begin{codeannotations}
\item 返回一个 \texttt{HandlerFunc}
\end{codeannotations}

内容目前一样。现在可以更新路由器了。

\begin{gocode}[title={代码清单 8.29 \texttt{router.go}：使用仓库}]
// The provided router is ready to serve.
func NewRouter(db *repository.GameRepository) *http.ServeMux {
    r := http.NewServeMux()
    // Register each endpoint.
    r.HandleFunc(http.MethodPost+" "+api.NewGameRoute,
        newgame.Handler(db)) #1
    r.HandleFunc(http.MethodGet+" "+api.GetStatusRoute,
        getstatus.Handler(db))
    r.HandleFunc(http.MethodPut+" "+api.GuessRoute, guess.Handler(db))
    return r
}

#1 Passes the dependency as a parameter
\end{gocode}

\begin{codeannotations}
\item 把依赖作为参数传入
\end{codeannotations}

怎么测试呢？我们只能把一个具体的仓库传给 \texttt{Handler} 函数。一旦用了真实的外部数据库，为了跑单元测试就得启动一个实例并连接上去。这绝对不可持续。我们用一个接口把它抽象出来。

\texttt{NewGame} 端点只需要向仓库添加一局游戏，别的什么都不需要。我们实际上可以通过定义一个最小接口来防止它做别的事。

% SOURCE_PDF_PAGE: 105
\begin{gocode}[title={代码清单 8.30 \texttt{newgame/handler.go}：最小接口}]
type gameAdder interface { #1
    Add(game session.Game) error
}

// Handler returns the handler for the game creation endpoint.
func Handler(db gameAdder) http.HandlerFunc { #2
    return func(w http.ResponseWriter, _ *http.Request) {
        // ...
    }
}

#1 gameAdder can only add a game.
#2 Uses the mockable interface
\end{gocode}

\begin{codeannotations}
\item \texttt{gameAdder} 只能添加游戏
\item 使用可 mock 的接口
\end{codeannotations}

路由器的 \texttt{db} 变量自动实现了这个小接口，现在在单元测试中创建一个只有一个方法的存根（stub）就很容易了。适配测试只需要两处修改。

\begin{gocode}[title={代码清单 8.31 \texttt{newgame/handler\_test.go}：用存根替代仓库}]
func TestHandle(t *testing.T) {
    handleFunc := Handler(gameAdderStub{}) #1
    req, err := ...
    // ...
    handleFunc(recorder, req) #2
    assert...
}

type gameAdderStub struct { #3
    err error
}

func (g gameAdderStub) Add(_ session.Game) error {
    return g.err
}

#1 Creates an anonymous function
#2 Calls it
#3 Stubs the repository
\end{gocode}

\begin{codeannotations}
\item 创建匿名函数
\item 调用它
\item 用存根替代仓库
\end{codeannotations}

你可以把这套逻辑推广到另外两个端点。GetStatus 只需要一个查找器，Guess 需要调用两个方法。现在，先别急着欢庆，有一件事要注意：还记得我们的服务器在不同 goroutine 中接收请求，

% SOURCE_PDF_PAGE: 106
并且并发地处理它们吗？向映射写入不是线程安全的。这意味着如果两个不同的例程向同一个映射写入，我们无法保证谁会胜出——甚至无法保证有谁能胜出。要解决这个问题，可以用上一章见过的概念：互斥锁。目标是避免对映射的并发访问，确保我们的服务器神智健全。

\subsection{给仓库添加互斥锁}

动机与上一章相同：我们要保护仓库——这里是那个映射——免受并发访问。保持简单，在 \texttt{internal/repository/memory.go} 的 \texttt{Add()}、\texttt{Find()} 和 \texttt{Update()} 方法里使用 \texttt{sync.Mutex}。

\begin{infobox}{注意}
互斥锁应当放在离线程不安全变量尽量近的地方。理论上我们可以在服务器上放一把互斥锁，一次只处理一个请求——但那会彻底违背微服务的初衷。想象一下，你最喜欢的网站一次只允许一个人访问！另一方面，如果存储方案本身已经线程安全——比如一个一次只接受一条查询的数据库——那我们根本不需要互斥锁。
\end{infobox}

首先，把互斥锁加到想保护的资源——\texttt{storage} 映射——旁边，放进 \texttt{GameRepository} 结构体。下面的清单给出代码。

\begin{gocode}[title={代码清单 8.32 \texttt{memory.go}：给 \texttt{GameRepository} 添加互斥锁}]
// GameRepository holds all the current games.
type GameRepository struct {
    mutex     sync.Mutex #1
    storage   map[session.GameID]session.Game
}

#1 Adds mutex to the storage
\end{gocode}

\begin{codeannotations}
\item 给存储添加互斥锁
\end{codeannotations}

% SOURCE_PDF_PAGE: 107
然后，每个方法都能通过接收者访问这个互斥锁。这是 \texttt{Add} 方法的示例代码。

\begin{gocode}[title={代码清单 8.33 \texttt{memory.go}：在 \texttt{Add()} 上调用 \texttt{Lock} 与 \texttt{Unlock} 互斥锁}]
// Add inserts a game in memory the first time it is received.
func (gr *GameRepository) Add(game session.Game) error {
    log.Print("Adding a game...")
    // Lock the reading and the writing of the game.
    gr.mutex.Lock() #1
    defer gr.mutex.Unlock() #2
    _, ok := gr.storage[game.ID]
    if ok {
        return fmt.Errorf("%w (%s)", ErrConflictingID, game.ID)
    }
    gr.storage[game.ID] = game
    return nil
}

#1 Locks the mutex before executing any read or write operation
#2 Executes the unlocking at the end of the method in a defer
statement
\end{gocode}

\begin{codeannotations}
\item 在执行任何读写操作之前锁定互斥锁
\item 在方法结束时以 defer 语句执行解锁
\end{codeannotations}

现在你可以自己更新其他方法，然后跑测试。如果测试通过，你的代码可以安全提交——开玩吧！

\section{改造 Gordle 库}

从第 5 章、或从我们的仓库复制 \texttt{gordle} 库。它远比我们需要的复杂，因为它包揽了整个会话。我们需要按此处的需求重构并简化这个库。我们知道需要以下用例：

\begin{itemize}
  \item 创建并返回一局新的 Gordle 游戏。
  \item 接受一次猜测并返回反馈。
\end{itemize}

我们本可以让 Gordle 库负责玩家剩余的尝试次数，但我们决定
% SOURCE_PDF_PAGE: 108
把这一逻辑放进会话（session）里。来分一分职责，如表 8.2 所示。

\medskip
\noindent\textbf{表 8.2 包的职责}

\begin{center}
\small
\begin{tabularx}{\textwidth}{|X|X|}
\hline
\textbf{\texttt{session} 包的游戏} & \textbf{Gordle 包的游戏} \\
\hline
\begin{itemize}[leftmargin=1.5em,itemsep=0.1em,topsep=0.2em]
  \item 存储剩余尝试次数
  \item 存储尝试与反馈的列表
  \item 通过调用 Gordle 来进行一次猜测
\end{itemize}
&
\begin{itemize}[leftmargin=1.5em,itemsep=0.1em,topsep=0.2em]
  \item 使用词库（corpus）为其答案创建一局游戏
  \item 接受一次猜测并计算反馈
  \item 允许通过反馈判断游戏是否结束
\end{itemize}
\\
\hline
\end{tabularx}
\end{center}

现在需要把 \texttt{session} 与 \texttt{gordle} 包的关注点分开。\texttt{gordle} 库不需要 ID，但会话需要。库不需要之前的猜测，但必须能告知游戏的状态。

\subsection{库的 API}

考虑到我们之前关于职责分配的决定，我们希望能够创建游戏并进行游戏。以下是我们需要暴露的东西：

\begin{gocode}[]
var g gordle.Game = gordle.NewGame(corpus)
var feedback gordle.Feedback = g.Play(guess)

var solution string = g.ShowAnswer()
var won bool = feedback.GameWon()
\end{gocode}

这应当够了。给反馈加一个 \texttt{fmt.Stringer}，就万事俱备。不过，还有几个设计问题要回答：

\begin{itemize}
  \item 谁负责告诉我们允许多少次尝试、还剩多少次？
  % SOURCE_PDF_PAGE: 109
  \item 谁负责读取词库文件并从中随机挑选一个单词？
\end{itemize}

这可以讨论，有几个选项可供选择。第一个选项：库把答案作为字符串参数接收，因为它不需要别的。第二个选项：可以说任何对该库的使用都需要这个行为，那不如就把它暴露出来。第三个选项：创建一个独立的、可测试的词库读取包。

在我们的例子里，由于大部分逻辑是从前一章复制来的，我们选择了第二个选项，让旧代码保持原样，直到我们证明它应当被进一步重构。因此，我们暴露一个返回字符串列表的 \texttt{ReadCorpus} 函数，而 \texttt{NewGame} 函数接收它并随机选择一个单词。注意，\texttt{New} 接收单一答案简化了测试，因为我们不必经过随机化。

我们不想让这本书被前一章复制粘贴来的大量代码拖垮，所以你可以在我们的仓库里找到简化后的最终包，或者自己动手玩玩、看看你需要什么。亲自体验一下给代码``瘦身''的乐趣吧。为了保持连贯、确保我们用的是同一份代码基础，我们包的最终 API 见下面的代码清单。

% SOURCE_PDF_PAGE: 110
\begin{gocode}[title={代码清单 8.34 包的 API}]
> go doc internal/gordle
package gordle // import "learngo-pockets/httpgordle/internal/gordle"
const ErrInaccessibleCorpus =
    corpusError("corpus can't be opened") #1
const ErrInvalidGuess = gameError("invalid guess length") #2
func ReadCorpus(path string) ([]string, error) #1
type Feedback []hint
type Game struct{ ... }
    func New(corpus []string) (*Game, error) #2

#1 Reads the corpus
#2 The game itself
\end{gocode}

\begin{codeannotations}
\item 读取词库
\item 游戏本身
\end{codeannotations}

阅读文档可以看出，第三个选项——把词库读取器放到别处——本可以让这个 API 更容易理解。欢迎按这个思路重构。如果你对库的测试覆盖率满意，是时候提交并使用它了。

\subsection{在端点中的使用}

到这里，我们已拥有编写端点逻辑所需的一切。我们已经在这三个端点里各自创建了一个函数来隔离那部分逻辑，所以只需在每处填充这个函数即可。

第一件事是把一个 \texttt{gordle.Game} 保留在我们的领域里。如果给 \texttt{session.Game} 加一个这种类型的字段，就能用这个字段来玩游戏，如下面的代码清单所示。有了这个新类型（并且测试通过），就可以完善端点了。

% SOURCE_PDF_PAGE: 111
\begin{gocode}[title={代码清单 8.35 \texttt{session/game.go}：在会话中使用 \texttt{gordle.Game}}]
// Game contains the information about a game.
type Game struct {
    ID           GameID
    Gordle       gordle.Game #1
    AttemptsLeft byte
    Guesses      []Guess
    Status       Status
}

#1 Adds a new field for the game
\end{gocode}

\begin{codeannotations}
\item 为游戏添加一个新字段
\end{codeannotations}

\subsubsection*{NEWGAME}

这里我们需要创建一局游戏，为它生成并保存一个随机 ID，然后返回该 ID。为什么要随机？递增 ID 是一个可怕的安全漏洞：任何人都可以创建一局游戏，然后到处乱玩、干扰别人的游戏（注意，身份验证将在 8.6.2 节讨论）。

我们可以用整数。如前所述，生成随机整数用 \texttt{math/rand} 包就相当快，但随机性较差；用 \texttt{crypto/rand} 包分布更好，但成本更高。

还有其他替代品，比如通用唯一标识符（UUID）——Go 中通常使用 Google 的库——或者按字典序可排序的通用唯一标识符（ULID）。我们选了最后这个，将使用 \url{https://github.com/oklog/ulid} 上的生成库。

如果你使用词库的相对路径，那么它需要相对于编译后的二进制文件来定义——不是定义它的那个文件，不一定是 \texttt{main.go}，也不是执行路径。

% SOURCE_PDF_PAGE: 112
\subsubsection*{使用 embed 包}

Go 标准库带有非常特别的 \texttt{embed} 包。这个包实现了一个极其有用的特性：它允许把普通文件的内容嵌入我们程序的源代码。为此，我们使用一种在编译时发给编译器的特定指令，称为指示（directive）或编译指示（pragma）。它们以注释的形式出现，形如 \texttt{//go:\{something\}}。

指示相当常见，用于向代码分析工具（如 linter）指出代码的特定部分。我们需要的语法如下：

\begin{gocode}[]
//go:embed corpus/english.txt
var englishCorpus string
\end{gocode}

这个 \texttt{englishCorpus} 变量必须定义在任何函数之外（放在函数里会被编译器当成普通注释——然后被忽略）。要用它，首先需要导入 \texttt{embed} 包。然而，因为我们不会使用该包的任何常量、变量、类型或函数，你的 IDE 可能会把这一行完全删掉（其实也不算冤枉——毕竟，不用的包为什么要导入呢？）。所幸，有个技巧可以强制导入包：把它别名为 \texttt{\_}（下划线）。这样包就不会被丢弃，会在导入它的地方可用。

以 \texttt{\_} 导入包通常发生在需要调用某些库的 init 函数时。这里，它确保 \texttt{embed} 包正确加载位于 \texttt{corpus/english.txt} 路径的文件内容——该位置相对于源 \texttt{.go} 文件。现在，文件内容已被加载进一个变量。

% SOURCE_PDF_PAGE: 113
\begin{gocode}[title={代码清单 8.36 \texttt{newgame/handler.go}：端点逻辑}]
func createGame(db gameAdder) (session.Game, error) {
    corpus, err := gordle.ReadCorpus(
        "corpus/english.txt") #1
    if err != nil {
        return session.Game{}, fmt.Errorf("unable to read corpus: %w", err)
    }
    game, err := gordle.New(corpus)
    if err != nil {
        return session.Game{},
        fmt.Errorf("failed to create a new gordle game")
    }
    g := session.Game{
        ID:           session.GameID(
            ulid.Make().String()), #2
        Gordle:       *game,
        AttemptsLeft: maxAttempts, #3
        Guesses:      []session.Guess{},
        Status:       session.StatusPlaying,
    }
    err = db.Add(g) #4
    if err != nil {
        return session.Game{}, fmt.Errorf("failed to save the new game")
    }
    return g, nil
}

#1 This path should come from the configuration.
#2 Generates the ID
#3 Uses a constant
#4 Saves the game for future plays
\end{gocode}

\begin{codeannotations}
\item 这个路径应当来自配置
\item 生成 ID
\item 使用一个常量
\item 保存游戏，供以后玩
\end{codeannotations}

我们选择用常量表示最大尝试次数；但如果你的词库包含不同长度的单词，你尽可以更有创意：从单词长度、难度设置、玩家的 ELO 或任何其他变量推导这个最大值。

这里，每次调用新游戏端点都会读一遍词库，看起来是在浪费资源。理想情况是在启动时加载它、在那时处理任何错误（比如 File Not Found），如果什么都没有就无法启动。如果服务连一份单词列表都拿不到，那它干脆别启动了。

% SOURCE_PDF_PAGE: 114
现在我们能拿到答案了，顺带也可以更新 API 适配器，补上 \texttt{WordLength} 和其他可能一直缺着的字段。此外，这个写死的词库路径一旦开始测试就会变成麻烦。

\subsubsection*{用正则表达式测试 NEWGAME 处理器}

怎么测试我们的处理器？只有一个坑：随机化。我们没法指望这个 \texttt{createGame} 函数被多次调用时输出完全一致。

在这种情况下，重要的是确定我们想测什么。我们可以用一个小技巧：把一个 ID 生成器接口作为参数传入，测试时提供一个返回固定值的生成器，现实中用返回 ULID 的生成器。我们也可以约定：需要断言的只是``ID 由那么多个字母数字字符组成''；满足这一点就满意了。至于``连续调用是否返回不同 ID''归不归我们管，也有争论空间。我们可以信任库做到了，但能信任我们永远记得调用库吗？

要验证 ID 由那么多字母数字字符组成，诀窍是使用正则表达式。翻翻 \texttt{regexp} 包，你会发现一大堆选项。

测试处理器本身时，我们可以把 ID 替换成已知字符串，如代码清单 8.37 所示。为此，我们把生成的 ID 从 JSON 输出中隔离出来，用 \texttt{strings.Replace} 替换它——该函数把旧字符串的给定数量的实例替换为新的（最后一个参数用 \texttt{-1} 表示替换所有出现）。

% SOURCE_PDF_PAGE: 115
\begin{gocode}[title={代码清单 8.37 \texttt{newgame/handler\_test.go}：在 JSON 输出中替换 ID}]
// idFinderRegexp is a regular expression that will ensure the body
// contains an id field with a value that contains
// only letters (uppercase and/or lowercase) and/or digits.
idFinderRegexp := regexp.MustCompile(
    `.+"id":"([a-zA-Z0-9]+)".+`) #1
id := idFinderRegexp.FindStringSubmatch(body) #2
if len(id) != 2 { #3
    t.Fatal("cannot find one id in the json output")
}
body = strings.Replace(body, id[1], "123456", 1) #4
assert.JSONEq(t, testCase.wantBody, body) #5

#1 Regexp to locate the ID
#2 Finds match in the JSON body
#3 Checks the number of occurrences
#4 Replaces the ID with a known string
#5 Compares the content of a JSON easily
\end{gocode}

\begin{codeannotations}
\item 用于定位 ID 的正则表达式
\item 在 JSON 体中查找匹配
\item 检查出现次数
\item 把 ID 替换为已知字符串
\item 轻松比较 JSON 内容
\end{codeannotations}

我们需要在输出体中寻找 \texttt{"id":"<若干数字和字母>"}，所以要匹配的正则表达式是 \texttt{`.+"id":"([a-zA-Z0-9]+)".+`}。拆解一下这个正则表达式：

\begin{description}
  \item[\texttt{.+}] 表示一列一个或多个（\texttt{+}）任意（\texttt{.}）字符。
  \item[\texttt{"id":"}] 正是由一个双引号、字母 i、字母 d、一个双引号、一个冒号和一个双引号组成的字符串。
  \item[\texttt{([a-zA-Z0-9]+)}] 这是一个捕获块——圆括号代表的就是它。正则表达式并不匹配这对圆括号本身，它只匹配一个或多个（\texttt{+}）任意字母或数字（范围在 a 到 z、A 到 Z 或 0 到 9 的字符）。不过，圆括号会捕获它们匹配到的内容，由 \texttt{FindStringSubmatch} 返回。
  \item[\texttt{".+}] 这个字符串以一个双引号开头，后随一个或多个（\texttt{+}）任意（\texttt{.}）字符。
\end{description}

% SOURCE_PDF_PAGE: 116
这个字符串被反引号 \texttt{`} 包裹，这样我们就不必在 Go 字符串中转义双引号。期望的响应体包含已知字符串，现在就可以用 \texttt{assert.JSONEq} 或等价物了。如果 ID 不符合期望的格式（字母数字字符），\texttt{FindStringSubmatch} 将找不到它，只返回一项：完整字符串。测试就会失败。

现在，测试 create 函数本身时，我们拿不到已编码的响应体。可以改用 \texttt{FindStringIndex}，它找出正则表达式在字符串中最左匹配的位置。索引在那里，就说明没问题。testify 的 \texttt{Regexp} 函数封装了这个用法。

\begin{gocode}[title={代码清单 8.38 \texttt{newgame/handler\_test.go}：用正则表达式校验 ID}]
func Test_createGame(t *testing.T) {
    corpusPath = "testdata/corpus.txt"
    g, err := createGame(gameCreatorStub{nil}) #1
    require.NoError(t, err)
    assert.Regexp(t, "[A-Z0-9]+", g.ID) #2
    assert.Equal(t, uint8(5), g.AttemptsLeft) #3
    assert.Equal(t, 0, len(g.Guesses))  #3
}

#1 Declares a stub that returns no error
#2 Matches the ID against a regexp
#3 Validates other fields ad lib
\end{gocode}

\begin{codeannotations}
\item 声明一个不返回错误的存根
\item 把 ID 与正则表达式匹配
\item 随意校验其他字段
\end{codeannotations}

正则表达式威力巨大，但在你不知道自己在看什么时也极其难懂。每当你写了一个，不要指望下一位维护者（包括未来的你自己）会觉得它容易解析：请系统地加一条注释，告诉他们它在找什么。确认测试通过、覆盖到位，就可以进入第二个端点了。

% SOURCE_PDF_PAGE: 117
\subsubsection*{GETSTATUS}

这里我们只需要调用数据库并返回游戏。就这样。当然，还要处理各种错误。呃，怎么处理错误呢？我们怎么知道是游戏没找到，还是出了其他意外错误（比如，在真实数据库情形下的连接错误）？

我们希望在游戏不存在时返回 Not Found 状态，其他情况返回 Internal Error。所幸，仓库暴露了一个我们可以对照检查的特定错误。

\begin{gocode}[title={代码清单 8.39 \texttt{getstatus/handler.go}：处理错误}]
game, err := db.Find(session.GameID(id)) #1
if err != nil {
    if errors.Is(err, repository.ErrNotFound) { #2
        http.Error(w, "this game does not exist", http.StatusNotFound)
        return
    }
    log.Printf("cannot fetch game %s: %s", id, err)
    http.Error(w, "failed to fetch game", http.StatusInternalServerError)
    return
}

#1 Checks in the repository for the game
#2 Checks against a specific error
\end{gocode}

\begin{codeannotations}
\item 在仓库中查找游戏
\item 对照特定错误检查
\end{codeannotations}

注意，这里我们选择不让错误``冒泡''到响应里：\texttt{http.Error} 的消息不包含 \texttt{err} 的值。把服务的内部细节暴露给客户端，很少是个好主意。随错误一起发回的措辞掩盖了真实错误的细节，所以我们需要把它记入日志以便调试。

\subsubsection*{GUESS}

最后，来玩我们的游戏！这里我们需要取回游戏、玩出这个单词、保存结果并返回它。

% SOURCE_PDF_PAGE: 118
\subsubsection*{可能的错误}

可能出什么岔子？成问题的场景可能包括：游戏没找到、存储没有响应、提交的单词不合法，以及游戏已经结束——或输或赢。还有一件事我们没加：领域（\texttt{session} 包）中的哨兵错误（sentinel error），用来告诉我们：不，你不能玩一局你已经赢了（或输了）的游戏。

先看看这个包办一切工作的函数必须做什么。我们先省略错误，然后再逐一思考每种情况。

% SOURCE_PDF_PAGE: 119
\begin{gocode}[title={代码清单 8.40 \texttt{guess/handler.go}：端点逻辑}]
func guess(id session.GameID, guess string, db gameGuesser)
    (session.Game, error) {
    game, err := db.Find(id) #1
    if game.AttemptsLeft == 0 || #2
        game.Status == session.StatusWon {
        return session.Game{}, session.ErrGameOver
    }
    feedback, err := game.Gordle.Play(guess) #3
    game.Guesses = append(game.Guesses, session.Guess{ #4
        Word:     guess,
        Feedback: feedback.String(),
    })
    game.AttemptsLeft -= 1 #5
    switch { #6
    case feedback.GameWon():
        game.Status = session.StatusWon
    case game.AttemptsLeft == 0:
        game.Status = session.StatusLost
    default:
        game.Status = session.StatusPlaying
    }
    err = db.Update(game) #7
    return game, nil
}

#1 Does the game exist?
#2 Are attempts still allowed?
#3 How does this guess fare?
#4 Records the play
#5 Nearing the end of the game
#6 Updates the status
#7 Saves the game for the next round
\end{gocode}

\begin{codeannotations}
\item 游戏存在吗？
\item 还允许尝试吗？
\item 这次猜得如何？
\item 记录这次游玩
\item 逐渐接近游戏尾声
\item 更新状态
\item 保存游戏，供下一轮使用
\end{codeannotations}

这是个很长的函数。每种情况下，错误应当是什么？

查找游戏时，我们可以简单地用一些上下文包装错误。这里我们用的上下文不多，但算个示例。这个错误的类型将永远是 \texttt{repository.Error}，而 \texttt{repository.ErrNotFound} 就可能从这里返回：

% SOURCE_PDF_PAGE: 120
\begin{gocode}[]
game, err := db.Find(id)
if err != nil {
    return session.Game{}, fmt.Errorf("unable to find game: %w", err)
}
\end{gocode}

玩游戏时也可能出错，那会是一个 \texttt{gordle.Error}：

\begin{gocode}[]
feedback, err := game.Gordle.Play(guess)
if err != nil {
    return session.Game{}, fmt.Errorf("unable to play move: %w", err)
}
\end{gocode}

最后，我们再次调用存储：

\begin{gocode}[]
err = db.Update(game)
if err != nil {
    return session.Game{}, fmt.Errorf("unable to save game: %w", err)
}
\end{gocode}

因为错误就是值，我们收到错误时就知道发生了什么，从而可以调整 HTTP 响应的状态码和消息。

\begin{gocode}[title={代码清单 8.41 \texttt{guess/handler.go}：使用 \texttt{domain.Game}}]
game, err := guess(id, r, db)
if err != nil {
    switch {
    case errors.Is(err, repository.ErrNotFound):
        http.Error(w, err.Error(), http.StatusNotFound)
    case errors.Is(err, gordle.ErrInvalidGuess):
        http.Error(w, err.Error(), http.StatusBadRequest)
    case errors.Is(err, session.ErrGameOver):
        http.Error(w, err.Error(), http.StatusForbidden)
    default:
        http.Error(w, err.Error(), http.StatusInternalServerError)
    }
    return
}
\end{gocode}

别忘了测试。没有什么窍门，但边界情况一大堆。

你拥有一个能运转的服务了！恭喜。现在想写一个前端客户端吗？只用 \texttt{curl} 玩游戏实在谈不上用户友好。你甚至可以用 Go 写一个
% SOURCE_PDF_PAGE: 121
调用该服务的 CLI。不过在放你走之前，我们还得考虑一下我们所走捷径的几点警告。

\section{安全概念与改进}

编写服务器时，我们要记住：目标终究是把它部署到某处去服务请求。然而，我们并不总知道一台服务器要负责多少用户、他们会发来多少查询。谈到安全，我们必须绞尽脑汁，想出所有``不高兴''的路径。我们需要识别什么可能出岔子，以防止最坏情形发生。在这里，有时不妨把自己想象成一心想找漏洞、搞垮系统的人。

\subsection{限制同时服务的请求数量}

针对服务器最常见的攻击之一叫作分布式拒绝服务（DDoS）攻击——其过程就是用过多请求把服务器压垮。这类攻击不会从服务器中提取任何信息，但会让服务器崩溃，使它对其他用户不可用。这种攻击通常让大量计算机向一台服务器发送成千上万的请求。每个请求都会使服务器分配内存去处理——一个并行任务（线程或 goroutine）、一些栈分配，等等。由于服务器的资源有限，到了某个临界点，海量的请求会让这些资源枯竭，服务器轻则什么都处理不了。

% SOURCE_PDF_PAGE: 122
前面已经看到，命中一个端点会启动一个 goroutine 来服务我们的请求。所幸 goroutine 很轻，成千上万个可以在单个线程里运行——但请求过载的风险始终存在。一些开源 mux 库提供节流（throttling）功能——限制 mux 在给定时刻能同时服务的请求数量。这是（尚）未进入 Go 标准库的特性之一。

\subsection{用户身份验证}

我们提到过，调用服务的客户端应当经过身份验证，这样我们才能确保一个玩家创建的游戏只有这一个人能玩。这也有助于限制每个用户的请求数，从而限制服务的负载。

怎么验证用户身份？一个非常常见的协议是开放授权（Open Authorization，OAuth）——通常写作其最新版本 OAuth 2.0。它用于通过一个身份验证服务器验证用户。典型做法是，验证用户是某一个服务的职责，它处理所有安全事务并提供一个签名令牌（token）。我们的 Gordle 服务随后通过 HTTP 头接收这个令牌，验证并解码它，然后在令牌中找到用户标识符。

根据需求，你可以决定客户端的身份验证级别：可以要求每个玩家都提供身份，或者对每个连接到服务的应用做验证。在第二种方案里，一个网站和一个移动应用会拥有不同的 ID 和密钥，各自有请求速率上限，而不管它们服务多少玩家。

身份验证及其解决的安全问题本身可以写满一整章，但它超出了
% SOURCE_PDF_PAGE: 123
本书的范围。想了解更多身份验证的内容，请阅读 Justin Richer 与 Antonio Sanso 的《OAuth2 in Action》（Manning，2017，\url{www.manning.com/books/oauth-2-in-action}）。

\subsection{日志}

本章我们用的是原生且非常基础的 \texttt{log} 包。在生产环境里这是个糟糕的主意：在并发环境——典型如服务——中，它会把日志输出搅得一团糟。市面上有很多优秀的日志库，能保护你的输出并提供格式化选项。Go 1.21 在标准库中加入了结构化日志器 \texttt{slog} 包——我们建议用它取代 \texttt{log} 包。

\subsection{错误格式化}

我们返回错误时，它就是一个简单的字符串。良好实践告诉我们：当端点的输出以 JSON 格式化时，错误也应当被格式化。比如这样：

\begin{gocode}[]
{
  "error": "game over"
}
\end{gocode}

这样，无论响应的状态码是什么，你的客户端都能使用同一个解码器。

\subsection{解码查询参数}

前面见过，查询参数在 API 中用作可选参数。它们总是成对出现：一个键和一个值。它们用作筛选器，指定要创建、更新或删除的资源。

% SOURCE_PDF_PAGE: 124
\subsubsection*{语法}

查询参数附加在路径之后，用一个问号与路径分隔。使用时，先写键，后跟等号，再写用来筛选的关联值。如果有多个参数，在每对之间加一个和号（\&）。问号之后的整个键=\texttt{值}对列表构成查询字符串（query string）。下一小节给出一个例子。

\subsubsection*{来实现一个例子}

我们想增加选择玩 Gordle 所用语言的可能。为此，把语言作为查询参数加入。创建游戏的 URL 将长这样：

\begin{gocode}[]
http://localhost:8080/game?lang=en
\end{gocode}

本例中，键是 \texttt{lang}，值是 \texttt{en}。先定义键的常量，并把它声明在路径参数常量旁边。

\begin{gocode}[title={代码清单 8.42 \texttt{internal/api/http.go}：作为常量的查询参数键}]
const (
    // GameID is the name of the field that stores the game's identifier
    GameID = "id"
    // Lang is the language in which Gordle is played.
    Lang   = "lang" #1
    // ...
)

#1 Defines the key for the query parameter
\end{gocode}

\begin{codeannotations}
\item 定义查询参数的键
\end{codeannotations}

这是有趣又容易的部分。接下来看如何从请求中解码查询参数。所有查询参数都可以从 URL 中取回，多亏了 \texttt{url} 库的 \texttt{URL.Query()} 方法，我们可以在处理器里通过
% SOURCE_PDF_PAGE: 125
\texttt{request.URL} 字段访问它。查一下返回类型，你会看到它是一个 \texttt{url.Values}，它（除其他外）暴露一个 \texttt{Get(key string) string} 方法：

\begin{shellcode}[]
> go doc url.Values
// Values maps a string key to a list of values.
// It is typically used for query parameters and form values.
// Unlike in the http.Header map, the keys in a Values map
// are case-sensitive.
type Values map[string][]string
...
\end{shellcode}

那么，在 \texttt{internal/handlers/newgame/handler.go} 文件中，用 \texttt{Values} 上的 \texttt{Get} 方法取回语言吧。

\begin{gocode}[title={代码清单 8.43 \texttt{newgame/handler.go}：在处理器中读取查询参数}]
// Handler returns the handler for the game creation endpoint.
func Handler(db gameAdder) http.HandlerFunc {
    return func(w http.ResponseWriter, r *http.Request) { #1
        lang := r.URL.Query().Get(api.Lang) #2
        if len(lang) > 0 {
            // TODO create a game in the chosen language
            fmt.Println(lang) #3
        }

#1 Replaces blank operator with a param for the request
#2 Retrieves the query parameter
#3 You can print it to check that it's working by checking the server
log.
\end{gocode}

\begin{codeannotations}
\item 用请求参数替代空白操作符
\item 取回查询参数
\item 你可以打印它，然后查看服务器日志来确认它正常工作
\end{codeannotations}

恭喜，关于 REST API 以及如何解码各类参数，你现在已经几乎无所不知了！

\section*{小结}

\begin{itemize}
\item web 服务是一个持续监听某个端口的程序，它基于一组处理器知道该如何处理请求。
\item 我们可以用 \texttt{http.NewServeMux()} 创建一个多路复用器，根据请求被发送到的 URL 把请求路由到特定的处理器。否则，就使用开源库。

% SOURCE_PDF_PAGE: 126
\item 处理器的任务是填满一个 \texttt{http.ResponseWriter}。应当对它 \texttt{Write}，有时还要用 \texttt{WriteHeader} 设置头。\texttt{ResponseWriter.Write} 设置的默认状态码是 \texttt{http.StatusOK}。如果想返回别的代码，需要在开始用 \texttt{Write} 写出响应内容之前调用 \texttt{WriteHeader}。
\item 如果端点里发生了错误，处理器不应把这个错误原样返回。如果你真想知道哪里出了问题，日志是错误最后展示的地方。
\item 有些目录名在 Go 中有特定含义。\texttt{testdata/} 里的文件不会被 \texttt{go build} 或 \texttt{go run} 编译。\texttt{internal/} 里的文件不能被其他模块导入。\texttt{vendor/} 由于历史原因是个应当避开的目录名。
\item 我们为 API 定义的类型应当只在端点处理器中使用。其余时间应当使用领域的类型。领域——或模型——类型不应当在服务模块之外可见。通常它们藏在 internal 目录里。
\item 仓库（repository 或 repo）提供对数据的访问，数据可以存在物理数据库、内存中，或任何形式里。
\item 检查你的代码是否线程安全永远是个好主意。编写测试，必要时使用互斥锁。
\item \texttt{embed} 包可用于在编译时加载文件（或目录）的内容。这很有用，比如你想把 SQL 查询保存在 \texttt{.sql} 文件里，或者像我们这样，想加载一组写死的值。
\item 一些（甚至大多数）开源包使用语义化版本。若不作指定，Go 原生会使用包的最新 \texttt{v1.x.y} 版本。要强制使用
% SOURCE_PDF_PAGE: 127
\texttt{v4.m.n}，你应当显式使用 ``\texttt{go get path/to/package/v4}''，并在 \texttt{.go} 文件中使用 \texttt{import "path/to/package/v4"}。
\item 使用正则表达式是匹配模式的利器。你会发现在各种情况下都会用到它，包括校验随机化的值。由于正则表达式会很快变得没法读，别忘了在注释里解释你写的每一个正则表达式。
\item REST API，或称 RESTful API，是一组定义两个系统之间接口的约束。REST API 通过 HTTP 通信，可以用 JSON、HTML 甚至纯文本交换数据。
\item HTTP 状态码非常有用，可以把服务器端发生的事情精确地告知客户端。HTTP 状态码是 1xx 到 5xx 的三位数字，可以表示一切顺利（200）或资源未找到（大名鼎鼎的 404）。返回正确的状态码意味着客户端能更好地理解发生了什么。拿不准时，返回 500；彻底没辙时，返回 418。
\item 通配符可用于在路由中表示路径参数。以模式 \texttt{"/games/\{id\}"} 注册处理器函数后，处理器就能用 \texttt{request.PathValue("id")} 访问 \texttt{id} 的值。
\item 同一个模式里可以用多个通配符：\texttt{"/users/\{user\_id\}/games/\{game\_id\}"} 可用于访问特定用户的特定游戏。
\end{itemize}
